\chapter{怎样把刑法问题写成可检查的推理}
\formalstatus

\section{先明确：什么能够证明，什么需要解释}
刑法规范不是从数学公理中推导出来的。立法决定哪些行为应受处罚，司法解释补充适用规则，案件证据支持事实认定；数学只能帮助我们检查：在这些前提之下，结论是否真正跟得出来。因而本册把四类句子分开书写：法律前提、解释假设、事实输入、形式推论。只有最后一类可以在前提固定后称为证明。把“通常如此”改写成一个双向箭头，并不会使它成为充分必要条件。

例如，“已经造成死亡”只描述结果，不能单独推出故意杀人罪。至少还须审查行为与死亡之间的归责联系、行为人的主观认识与态度、责任条件以及正当化事由。把案件缩成一个结果数值，再代入“死亡即有罪”的函数，遗漏的正是法律判断中最重要的部分。

本册中的所有设例均为解释模型而编写的简化案件。设例写明某条件成立，是演算的输入；现实案件中，该条件必须有相应证据支持。将两者区分，才能避免拿题目的假定替代司法证明。

\section{原书导读中的一个推理需要收窄}
原书“法考是不是‘三无考试’”一文，以行政许可法第54条禁止指定教材的规定，说明法考没有官方教材。本册原文重排与逻辑核查分开：保留作者原话，并在这里指出，\textbf{“行政指定”与“官方组织编写或出版”并不是同一个谓词。}

\lawinput{行政许可法第54条规定有关考试应公布报名条件、报考办法、考试科目和考试大纲，并禁止组织强制性考前培训及指定教材或其他助考材料。\lawref{L9}{第54条}}
令$S(x)$为“有关机关将$x$指定为教材”，$O(x)$为“$x$由官方组织编写或出版”，$F(x)$为“该行为落入本条禁止范围”。从条文可提炼相关禁止关系$S(x)\Rightarrow F(x)$，但仅凭该式，不能推出$O(x)\Rightarrow F(x)$，因为还缺少$O(x)\Rightarrow S(x)$这一桥接前提。

严格反模型可取一个对象$a$，令$O(a)=1,S(a)=0,F(a)=0$。此时$S\Rightarrow F$仍为真，而$O\Rightarrow F$为假。这证明了所述推论在逻辑上不成立；它不证明任何具体出版活动必然合法，也不证明某套教材具有法定效力。评价具体活动仍须核对事实与其他适用规范。

更准确的表达应限定为：该条禁止行政指定考试教材；某份参考材料是否属于被禁止的指定行为，需要另行判断。原文关于考试成绩是“唯一有效评价标准”等表述，也应标注为作者的评价主张；不能不经论证就作为刑法解释或教材质量的公理。至于历年公开试题、答案的实际发布情况，属于需要逐年查证的事实，本册不依据上述逻辑分析替代核查。

\section{统一输入：人、行为、时间、对象与规范版本}
\begin{modeldef}[案件输入]
一个案件模型记为
\[
 \mathcal X=(P,A,O,T,F,E,N,I).
\]
其中，$P$为行为人集合，$A$为行为集合，$O$为行为对象集合，$T$为时间轴；$F$为待认定事实，$E$为证据材料，$N$为适用的规范文本及版本，$I$为明确选择的解释方案。它们不是八项犯罪构成要件，而是开展分析所需的信息分类。
\end{modeldef}

为了防止同名词混淆，写作$\operatorname{Intent}(p,a,r,t)$，表示人在特定时间就某一行为可能导致的某一结果所具有的故意。不能把“对取财有故意”直接替换成“对致人死亡有故意”。同理，$\operatorname{Possess}(p,o,t)$中的占有必须写明时间；在$ t_0 $代为保管，在$ t_1 $擅自处分，两时刻的占有及财产关系可能不同。

规范版本记作$N=(\text{名称},\text{条号},\text{公布日},\text{施行日},\text{适用条件})$。只保存一个条号不够：同一条号可能修订。书名中的“2027”不能替代法律施行日期。对金额标准还应加入地域和行为时间，写成$\theta=\theta(N,\text{地区},t)$。

\section{未知不等于否定：三值记录与举证要求}
\lawinput{刑事诉讼法对定罪证据提出法定要求，包括定罪量刑事实有证据证明、证据经过法定程序查证属实，以及综合证据排除合理怀疑。本册的符号不改变这些要求。\lawref{L5}{第55条}}

在整理题目或卷宗时，先允许事实取值
\[
 v(P)\in\{1,0,?\},
\]
分别表示已给定或已证明、已否定、尚不明确。这里的“?”是工作记录，不是新增法律上的证明标准。对合取条件，可用下表检查信息是否齐全。

\begin{center}
\begin{tabular}{c|ccc}
$\wedge$ & $1$ & $0$ & $? $\\\hline
$1$ & $1$ & $0$ & $? $\\
$0$ & $0$ & $0$ & $0$\\
$?$ & $? $ & $0$ & $? $\\
\end{tabular}
\end{center}

若故意杀人罪的某项必备条件取“?”，不能为了算出一个答案而把它改成$1$；也不能说事实世界中它已被证明为$0$。证据不足不能定罪，是程序上如何处理证明缺口的问题，不能解释成数学上已经证明被告没有实施任何行为。

\begin{modelprop}[合取条件的否定与缺口]
设模型$C=P_1\wedge\cdots\wedge P_n$。若有一个$P_j=0$，则$C=0$；若没有条件为$0$，但至少一个条件为“?”，则$C$在本工作表中仍记为“?”。
\end{modelprop}
\begin{proof}
按照上表，合取式中与$0$相合取的结果均为$0$。若只有$1$和“?”，每加入一个$1$不改变已有的“?”；至少出现一个“?”就不能把全部合取写成$1$。这只证明工作表的运算规则，不证明任何具体案件的事实。
\end{proof}

\section{严格推导需要的五步}
一个可复核的论证至少写出以下五步。
\begin{enumerate}
\item 写出结论的准确对象：谁、哪个行为、哪个罪名、哪一阶段。
\item 列出规范前提，并区分法条、司法解释与争议学说。
\item 把事实逐项对应到条件；不能以“案情严重”替代缺失的要件。
\item 说明推理规则：合取、析取、反证、分类讨论，或规范明确规定的例外。
\item 给出边界：改变哪个事实，结论会改变；哪个问题仍需证据或解释。
\end{enumerate}

例如，$A\Rightarrow B$与$B\Rightarrow A$不是同一句话。法律说“出现甲条件可以从轻”，不能倒推“判了轻刑就一定出现甲条件”；法律说“具备甲乙丙条件应当减轻”，也不能把“应当”偷换成可任意选择的“可以”。

\begin{modelprop}[增加必要条件不会扩大适用集合]
在同一论域中，令$S=\{x:P(x)\}$，$S'=\{x:P(x)\wedge Q(x)\}$。则$S'\subseteq S$。
\end{modelprop}
\begin{proof}
任取$x\in S'$，依定义同时有$P(x)$与$Q(x)$，所以有$P(x)$，从而$x\in S$。这证明了集合包含关系。至于法律是否确实要求$Q$，仍须单独查明，不能由该包含关系推出。
\end{proof}

这条简单命题特别适合检查未成年人责任、特殊主体和情节加重等问题：如果漏掉了一个必要条件，适用范围通常被错误扩大；如果擅自增加法无根据的条件，范围通常被错误缩小。

\section{学习时如何记忆与自检}
记忆单位采用“条件—结论—边界”三联卡。卡片正面只写问题，背面写规则和一个反例；不能把所有信息都印在正面再误以为已经记住。每个规则至少准备两道变式：一题只改变一个必要条件，另一题保持结果相同而改变原因。

建议使用一套可调整的复习安排：学习当天、次日、第3日、第7日、第14日、第30日闭卷检索。该安排是练习方案，不是关于记忆效果的实验结论。若某张卡连续答错，回到法条和条件表，不机械增加背诵次数。错误分为四类记录：漏条件、错对象、错时间、错规范版本。

\begin{center}
\begin{tabular}{p{.16\textwidth}p{.34\textwidth}p{.39\textwidth}}
\toprule
卡片位置 & 要写的内容 & 自查问题\\\midrule
规则 & 适用前提及规范后果 & “可以”和“应当”是否区分？\\
对象 & 人、行为、结果、罪名 & 是否把甲的故意移给乙？\\
时间 & 行为时、结果时、介入时 & 是否用事后获知代替事前明知？\\
反例 & 缺少一个关键条件的案件 & 去掉它后为什么不能推出原结论？\\
版本 & 条号、来源与生效信息 & 是否把旧金额标准当现行标准？\\\bottomrule
\end{tabular}
\end{center}

\chapter{总论：规范、构成与归责}

\section{第一讲：刑法论}
\subsection{解释是确定规则的适用范围}
设一个法律用语为$w$，解释方案$I$给出其可适用对象集合$\llbracket w\rrbracket_I$。解释问题是在文义、体系、规范目的等理由约束下确定该集合，而不是先选择希望处罚的人，再倒着调整集合。文义解释、体系解释与目的解释是论证方法，不能分别贴上“对”或“错”的标签就代替具体说明。

\lawinput{罪刑法定要求刑事评价有法律依据，法律没有明文规定为犯罪的，不得定罪处刑。刑法第4条、第5条分别规定平等适用与罪责刑相适应原则。\lawref{L1}{第3至5条}}

在模型中设置准入约束$\operatorname{Authorized}(N,I,k)$，表示罪名$k$及采用的解释有合法依据。只有该约束通过，才继续考察事实是否落入构成集合。社会危害性的判断不能独立创造一个没有法律依据的新罪名。

\caseinput{某新型行为明显给他人造成损失，但查明其不满足任何现有罪名构成。}
推导为：第一，损失成立只是事实输入；第二，逐项检索现行罪名而均不成立；第三，无法得到具有依据的犯罪构成；第四，不能以“应当受到惩罚”为由补造罪名。行政责任、民事责任是否成立，需要另行适用相应规范。

\subsection{时间效力和空间效力分别计算}
把行为时间$t_a$、裁判时间$t_j$和各规范生效时间分开。在新旧法比较中，先判断行为时是否为犯罪，再按照法定从旧兼从轻原则处理；“新法编号更新”不推出“新法总是优先”。\lawref{L1}{第12条} 比较的对象是法律规定下的具体法律后果，不能随意从旧法取一个有利条件、再从新法取另一个有利条件拼出第三部法律。

空间效力用管辖连接点集合记录：领域内、船舶或航空器、国籍、保护原则、普遍管辖等各有前提。不同连接点可能并存，不能把一个连接点成立误当成其他程序性限制都已消失。\lawref{L1}{第6至11条}

\modelbound{本节形式化的是规范选择程序；解释结论仍需要文本和法理支持。区分扩大解释与类推解释尤其不能只比较两个集合的大小。}
\memorycard{依次写出“依据—解释—地域—时间”。反例卡：新法较重为什么不能只因裁判发生在新法施行后就适用新法？}

\section{第二讲：犯罪构成}
\subsection{三段论的正确位置}
原书以改良版四要件体系组织讲授。本册保留其检索顺序，但将每一罪名的条件单独列全。设$C_k(x)$为罪名$k$的构成符合性，$J(x)$为违法阻却事由，$R_k(x)$为责任条件。用于学习的候选规则写成
\[
 C_k(x)\wedge\neg J(x)\wedge R_k(x)
 \ \Longrightarrow\ \operatorname{Candidate}_k(x),
\]
其中箭头只在适用规范、证据要求、犯罪形态等已经明确的前提下使用。$C_k$内部不能只写“客体、客观、主体、主观”四个标签，必须展开为具体事实条件。

\caseinput{甲为非法占有他人手机，在知道手机由乙占有的情况下，偷偷取走并建立稳定控制。设甲达到相应责任年龄、具备责任能力；入罪条件已满足，无阻却事由。}
第一步，确定对象是他人占有的财物。第二步，确定取走行为与占有转移。第三步，确认非法占有目的及故意。第四步，核对法定入罪路径。第五步，审查责任和违法阻却。本例得到盗窃罪候选结论；若另有竞合问题，再进入罪数评价。

\subsection{必要条件不能交换为充分条件}
若$C_k\Rightarrow P$，$P$是成立$C_k$的必要条件。知道$P$成立，还不能推出$C_k$。例如盗窃罪通常涉及非法占有目的，但具有非法占有目的并不足以确定是盗窃：诈骗、抢劫等也可能具有这一目的。

\begin{modelprop}[共同条件不能区分罪名]
设$C_1\Rightarrow P$且$C_2\Rightarrow P$，并存在$x_1\in C_1\setminus C_2$及$x_2\in C_2\setminus C_1$。只观察$P=1$，不能从逻辑上唯一确定案例属于$C_1$还是$C_2$。这里用$C_i$同时表示相应谓词的真值集合。
\end{modelprop}
\begin{proof}
两个给定案例在观察量$P$上都等于$1$，且$x_1$只属于两类中的$C_1$，$x_2$只属于两类中的$C_2$。仅依$P$作出的唯一分类必然把至少一个案例误分，因而需要能够区分取财机制等附加条件。
\end{proof}
\modelbound{四要件体系和阶层体系是不同分析组织方式。用某一套符号排出顺序，不证明其他体系错误；争议应比较具体前提与结论。}
\memorycard{把一个罪名写成完整条件表，再依次删除一项并寻找反例。能说出“少了什么所以不能推出”，才算理解三段论。}

\section{第三讲：犯罪客体}
\subsection{保护利益不等于行为对象}
令$O(x)$表示案件中直接受到作用的对象，$L_k$表示罪名$k$保护的法益。手机是对象，财产利益是可能受保护的法益；人体是对象，生命、身体健康、性自主或人身自由可能是不同法益。一个物理对象可以与多个法益有关，同一法益也可以通过不同对象受侵害。

构造一个二部关系$\mathcal R\subseteq O\times L$，记录“对象与受侵害利益”的联系。该关系用于提出候选罪名，不能单独得出罪名。定罪仍须回到分则构成。

\caseinput{甲毁坏他人用于紧急救援的设备。}
第一步，财物受损提示财产法益；第二步，设备用途可能提示公共安全或救援秩序；第三步，分别检索有关构成的对象、危险状态、主观条件；第四步，在这些条件查明前，不能仅以设备价格决定全部评价。该例说明对象价值与全部法益侵害不是同一个变量。

\subsection{怎样避免“法益万能论”}
在逻辑上，保护目的通常是解释约束之一，而不是足以自动产生犯罪的充分条件：
\[
 \operatorname{HarmTo}(L_k,x)
 \not\Rightarrow C_k(x).
\]
反例结构是：某行为确实造成利益损失，但缺少法律规定的行为方式、主观要素或入罪条件。此时要件缺失不能以抽象法益补足。

\modelbound{不同学说对某一罪保护法益的具体表述可能不同。本册仅把它作为解释与检索工具，不将某个学说名称写成无需证明的定理。}
\memorycard{同一对象画两条线：一条写“作用于什么”，另一条写“损害什么利益”。再问：还缺哪个具体要件？}

\section{第四讲：犯罪主体}
\subsection{自然人、特殊身份和单位分别建模}
对自然人，记$\operatorname{Person}(p)$、$\operatorname{Capacity}(p,t)$、$\operatorname{AgeEligible}_k(p,t)$。对要求特殊身份的罪名，另设$\operatorname{Status}_k(p,t)$。身份是特定构成中的条件，不是对个人全部行为的永久标签。

对单位犯罪，至少记录单位主体资格、单位意志或依法可归属于单位的决策、行为与单位的联系，以及分则是否规定该罪可由单位构成。刑法第30条、第31条提供单位犯罪及处罚的基本框架。\lawref{L1}{第30、31条}

\caseinput{公司负责人以公司名义实施取财行为，所得实际全部归其个人。}
第一步，名称中出现“公司”并不足以证明单位犯罪；第二步，核对有关罪名是否规定单位主体；第三步，查明决策、利益归属与行为归属；第四步，分别分析自然人责任与单位责任。不能从公章或职位一个事实直接推出单位犯罪。

\subsection{身份作为关系，不作为模糊印象}
在职务犯罪中，$\operatorname{Status}(p,t)$与$\operatorname{DutyLink}(p,a,t)$必须分别证明。某人具有国家工作人员身份，并不意味着其所有私人活动都利用了职务便利。反过来，特定法律规定可能按实际履行公务等条件认定主体，不能只凭劳动合同名称排除。

\modelbound{没有特殊身份者能否作为共犯参与某种身份犯，需要结合罪名结构和参与方式分析；“不是特殊主体”不自动等于“对该事件完全无刑事责任”。}
\memorycard{主体卡写三栏：“谁”“行为时的身份”“身份如何连接该行为”。单位卡增加“分则依据、决策、利益归属”，不得只背“单位名义”。}

\section{第五讲：客观要件一——行为}
\subsection{作为与不作为都需要具体化}
对作为，记录动作、对象、工具、场所、时间以及行为的社会意义。对不作为，不能把“没有做某事”直接等同犯罪；世界上没有做的事情无限多，法律只在特定条件下要求某人履行特定义务。

学习模型记为
\[
 O= D\wedge A\wedge M\wedge E\wedge Q,
\]
其中$D$为存在特定作为义务，$A$为具有履行能力，$M$为没有履行，$E$为在有关不真正不作为犯中具有与作为相当的构成意义，$Q$为该罪的其余条件。真正不作为犯由具体条文规定义务与行为结构；不真正不作为犯还需要论证为何能按结果犯评价。上述划分属于分析模型，具体义务来源与等价判断必须说明根据。

\caseinput{负有照护义务的甲，能够及时为婴儿提供食物，却为使其死亡而持续拒绝提供，最终发生可归责的死亡。}
第一步，写明监护照护关系而非仅写“甲在现场”。第二步，证明有能力获得并提供食物。第三步，确认持续不履行。第四步，将死亡归于不履行所实现的危险。第五步，分析杀人故意及责任。只有这些条件一一成立，才可进入不作为故意杀人的评价。

\subsection{反事实避免不能机械等同归责}
设$R(0)$为实际未履行时的结果，$R(1)$为假定履行时的结果。若$R(0)=1,R(1)=0$，可以说明履行可能对避免结果有决定意义，但反事实是否可靠需要医学、工程等证据。即使可靠，也不能跳过法定作为义务与主观条件。

反例：会游泳的路人和负有救生职责的人都在岸边，二人所负具体义务可能不同；不能仅因为二人身体能力相同，就推出刑事责任相同。另一个反例：负有义务但客观上根本无法实施有效救助，不能把无能力履行的情形与有能力而拒不履行混同。

\modelbound{作为义务的来源、先行行为的范围和等价性在边缘案例中存在解释问题。应把竞争观点分别列为$I_1,I_2$，不能通过删掉争议参数伪装成唯一答案。}
\memorycard{闭卷展开“义务—能力—不履行—等价—其余要件”。每次用一名普通路人与一名有特定义务者作对照。}

\section{第六讲：客观要件二——结果}
\subsection{结果层级与构成要求}
记$H$为实害结果，$D_c$为具体危险，$B$为条文规定的行为。不同犯罪可能要求$B\wedge H$、$B\wedge D_c$，或依法以特定行为的实施作为构成核心。不能预设所有犯罪都必须等到实害发生，也不能把所有危险都视为已经满足某个具体危险犯。

危险判断需要明确对象与条件：什么利益、在什么时间、处于何种具体危险。单纯写“有风险”信息太少。对抽象危险犯，仍须核对条文和解释要求的具体行为，不得将抽象危险理解为无条件入罪。

\subsection{结果加重犯的拆分}
以$B$表示基本犯罪，$H^+$表示加重结果，$C$表示归责联系，$G^+$表示对加重结果至少具有规范要求的罪过，$L$表示法律规定该结果加重结构。学习模板为
\[
 B\wedge H^+\wedge C\wedge G^+\wedge L.
\]
五项缺一不能直接得到该结果加重评价。特别是$B\wedge H^+$并不自动等于结果加重犯。\lawref{L1}{第14至16条及相应分则条文}

\caseinput{甲故意实施伤害，乙死亡。}
第一步，确认甲是伤害故意还是杀人故意；第二步，查明伤害行为与死亡的归责联系；第三步，审查死亡是否符合有关结果加重规定及罪过要求；第四步，区分故意伤害致死、故意杀人、只有基本伤害责任或其他评价。死亡相同不能消除主观与因果差异。

\modelbound{结果加重犯的未遂等问题可能涉及具体罪名与学说分歧。本节只给出拆分方法，不从“出现死亡”机械推出任何单一罪名。}
\memorycard{结果卡先问“本罪要求实害、具体危险还是特定行为”，再问“加重结果凭哪一条法律连接基本行为”。}

\section{第七讲：客观要件三——因果关系}
\subsection{先作事实说明，再作法律归责}
用有向图描述案件中的先后与机制：行为$A$，中间变化$M$，结果$R$；另有介入因素$Z$。图上的箭头必须有事实机制或证据支持，时间先后本身不够。教学模型可写为
\[
 M=f(A,U_M),\qquad R=g(M,Z,U_R),
\]
其中$U_M,U_R$表示背景因素。改变$A$并观察$R$的反事实，是分析工具，不是直接用于实际案件的自动算法。

\begin{center}
\begin{tikzpicture}[>=Stealth,node distance=12mm and 19mm,every node/.style={draw,rounded corners,align=center,inner sep=5pt}]
\node (a) {行为 $A$};
\node[right=of a] (m) {危险的发展 $M$};
\node[right=of m] (r) {结果 $R$};
\node[below=of m] (z) {介入因素 $Z$};
\draw[->] (a)--(m);\draw[->] (m)--(r);\draw[->] (z)--(r);
\end{tikzpicture}
\end{center}

是否将结果归责于行为，还须分析行为制造或增加的危险、结果实现方式和规范保护范围。本册将其作为明确的解释分析框架，并不声称刑法用一个普遍公式规定了所有因果案件。

\subsection{介入因素不是一出现就中断}
\caseinput{甲使乙受伤，乙送医后因通常救治中的一般风险死亡。另作变式：乙仅受轻伤，但住院时遭无关第三人独立蓄意杀害。}
在第一例，先查初始伤害的危险、治疗必要性、死亡机制与医疗介入作用；不能只因出现医生就宣布甲与死亡无关。在第二例，应把第三人的独立杀害作为新的机制，分别分析初始伤害是否仍对该死亡有法律上的归责意义。两个例子的区别在机制与风险实现，不在介入者的人数。

严格写法是：先列出至少两个竞争解释$H_1,H_2$，再检验证据能够排除哪个解释。若死亡原因仍无法查明，不得以“甲先动手”补足归责缺口。对可能存在择一因果、多重充分原因等复杂情形，简单的“如果没有甲，结果还会不会发生”也可能失效，需要更精细的事实模型。

\subsection{一个逻辑反例}
两人各自实施足以导致同一损害的行为，若把必要条件反事实作为唯一标准，可能错误得出两人都不是原因。因此，必要条件检验的便利不能被误写成普遍充分必要定理。该反例只说明一种模型有适用限制，不能直接给出所有共同作用案件的法律结论。

\memorycard{顺序记为“行为危险—发展机制—介入作用—结果是否实现该危险”。每张因果题必须画一条时间线，再写一条机制线，两条线不能互相替代。}

\section{第八讲：主观要件}
\subsection{故意、过失和意外事件的分解}
\lawinput{故意犯罪要求对危害结果具有明知，并希望或者放任；过失分为应当预见而因疏忽没有预见，以及已经预见但轻信能够避免。过失犯罪只有法律规定才负刑事责任；无故意或过失的意外事件依相应规定处理。\lawref{L1}{第14至16条}}

令$K$表示对特定结果的认识，$W$表示希望，$A$表示接受其发生，$V$表示确信可避免且其认识结构属于轻信，$D$表示应当预见且能够预见。
\begin{align*}
 \text{直接故意的教学特征}&: K\wedge W,\\
 \text{间接故意的教学特征}&: K\wedge\neg W\wedge A,\\
 \text{过于自信过失的教学特征}&: K\wedge V,\\
 \text{疏忽大意过失的教学特征}&: D\wedge\neg K.
\end{align*}
这里以$\neg W$明确排除希望该结果发生的直接故意情形，避免仅以“接受”一词使两个教学类别重叠。这些式子是把法律要素拆开，不是用几个心理分数直接诊断人的内心。主观状态要结合言语、行为方式、工具、危险程度、事前准备、事后行为等证据综合判断。某一事后表现不具有永远固定的证明力。

\caseinput{甲知道某行为可能造成乙死亡。情形一：甲认为“死了也无所谓”，仍实施；情形二：甲依赖某项具体措施，轻信死亡不会发生，仍实施。}
第一步，两例都具有对死亡可能性的认识；第二步，区分对死亡的态度；第三步，核对“具体措施”是否为真实认识及其事实依据；第四步，才讨论间接故意与过于自信过失。不能把“明知可能发生”单独当作区分标准。

\subsection{事实认识错误要同时保留现实与认识}
用$F$表示实际事实，用$\widehat F$表示行为人的认识。评价时必须成对比较$(F,\widehat F)$，不能只看其中一边。对象错误、打击错误、因果过程错误、抽象事实认识错误应分别分析；采用法定符合说或其他学说时，注明解释立场。

\caseinput{甲误把他人所有的伞当作自己的伞取走。}
第一步，实际对象属于他人；第二步，甲认识中的对象属于自己；第三步，如果这种认识真实存在，非法占有他人财物的故意可能缺失；第四步，不能仅因客观上拿走他人财物就认定故意盗窃。若甲知道伞不是自己的，仅以为“拿走无所谓”，则认识对象不同，不能直接套用本例。

\modelbound{主观要件依罪名分别判断。不能把对一个结果的故意转给另一个结果，也不能把情绪恶劣等同于特定犯罪故意。}
\memorycard{固定追问三句：“知道什么”“对哪个结果抱什么态度”“该罪是否处罚这种过失”。错误论题用两栏表写现实与认识。}

\chapter{总论：阻却事由、犯罪形态与参与关系}

\section{第九讲：排除违法事由}
\subsection{正当防卫的条件不能只保留结果}
\lawinput{正当防卫围绕正在进行的不法侵害展开；防卫过当要求明显超过必要限度并造成重大损害。特殊防卫另有适用条件。最高法、最高检、公安部指导意见强调，必须立足防卫时的具体情境，不能以事后冷静的精确标准苛求防卫人。\lawref{L1}{第20条}\lawref{L4}{第5至18点}}

定义$U$为存在不法侵害，$T$为正在进行，$O$为针对侵害人，$P$为保护合法权益的目的，$X$为明显超过必要限度，$H$为造成重大损害，$S$为满足特殊防卫条件。另令$A$表示不存在指导意见第10点等所述滥用防卫权而应排除防卫行为性质的情形。令基础防卫条件$B=U\wedge T\wedge O\wedge P\wedge A$。在本教学模型所概括的通常情形中：
\begin{align*}
 \text{一般防卫中不成立过当的分支}&:B\wedge\neg S\wedge\neg(X\wedge H),\\
 \text{防卫过当候选}&:B\wedge\neg S\wedge X\wedge H,\\
 \text{特殊防卫}&:B\wedge S.
\end{align*}
这里“候选”表示仍需根据主观罪过、具体罪名和责任条件完成犯罪评价。对借防卫之名实施侵害、防卫挑拨等情形，不能先认定$B$成立再用本表计算。

\begin{modelprop}[防卫过当的双条件结构]
若规范要求$X\wedge H$，那么只证明$H$，或只证明$X$，均不足以推出防卫过当。
\end{modelprop}
\begin{proof}
取$X=0,H=1$，则$X\wedge H=0$；取$X=1,H=0$，仍有$X\wedge H=0$。两组赋值分别提供反例，说明单项不足。法律上如何判定$X,H$须依据规范与案情，并非由这张真值表自动完成。
\end{proof}

\caseinput{乙持刀持续攻击甲的要害，甲为制止侵害反击，导致乙死亡。}
先确认现实严重侵害及其持续状态，再确认对象、目的与特殊防卫条件。若符合第20条第3款，不能仅因死亡结果反推防卫过当。变式：乙已确实放弃并失去继续侵害能力，甲随后报复；此时首先改变的是时间条件，不能沿用原防卫结论。

\subsection{紧急避险和被害人承诺}
紧急避险可用$D\wedge T\wedge N\wedge L$组织：现实危险、紧迫状态、不得已的必要性、合法限度，并另查主体例外与其他要求。利益比较不是给生命、尊严等随意标价，而是在法规范约束下判断是否超过必要限度。\lawref{L1}{第21条} 不能把避险中对第三人的损害直接套成防卫对侵害人的损害。

对被害人承诺，至少检索可处分利益、承诺能力、真实自愿、范围与时间、行为是否仍被特别规范禁止。承诺不是一个适用于所有罪名的万能免责变量。某人同意他人实施致命侵害，不因此自动使行为合法。

\memorycard{防卫先背“起因、时间、对象、目的”，再用两格检验“明显超限、重大损害”，最后单独检查特殊防卫。避险另开卡片，防止混用。}

\section{第十讲：排除责任事由}
\subsection{责任年龄是分段条件，不是单一阈值}
\lawinput{第17条经刑法修正案（十一）修改。满16周岁的一般责任规则，14至16周岁的特定严重犯罪规则，12至14周岁的极严格例外规则，各有不同条件；对依法追究责任的不满18周岁者，应当从轻或者减轻。\lawref{L2}{第1条}}

用$a$表示行为时已满周岁数，$K_8$表示14至16岁范围内法律列举的八类严重行为条件；用$K_{12}$表示故意杀人、故意伤害的限定犯罪类型，$D$表示死亡，$C$表示特别残忍手段致重伤并造成严重残疾，$V$表示情节恶劣，$A$表示最高人民检察院核准追诉。年龄准入模型为
\[
 E_a=
 \begin{cases}
 0,&a<12,\\
 K_{12}\wedge(D\vee C)\wedge V\wedge A,&12\le a<14,\\
 K_8,&14\le a<16,\\
 1,&a\ge16.
 \end{cases}
\]
$E_a=1$仅表示满足相应年龄门槛，不表示其他犯罪条件自动成立。八类条件为故意杀人、故意伤害致重伤或死亡、强奸、抢劫、贩卖毒品、放火、爆炸、投放危险物质；具体案件按法定范围分析，不能随意扩成“一切严重犯罪”。

\caseinput{13周岁的甲故意伤害乙造成重伤，但并未造成严重残疾，且没有死亡。}
即使行为严重，$(D\vee C)$仍未依给定事实成立，因此不能只凭“已满12岁”推出第17条第3款刑事责任。该结论仅针对所给低龄追诉条件，其他保护、矫治制度需另行判断。

\subsection{责任能力的“或者”不可误写为“并且”}
对精神障碍，刑法第18条的核心判断包括行为时不能辨认或者不能控制自己行为，并需经法定程序鉴定确认。\lawref{L1}{第18条} 把“不能辨认或者不能控制”误写成“两者都不能”，会缩小法定免责条件。间歇性精神病人在精神正常时犯罪、尚未完全丧失能力、醉酒等情形应分别按条文分析。

违法性认识可能性与期待可能性属于需要谨慎处理的责任理论问题。不能从“不知道具体条号”直接推出无责任；也不能自创“压力很大即免责”的规则。必须说明所采用理论、法律依据以及为何达到相应程度。

\memorycard{年龄题画四段数轴，并在12至14岁段下面写四重条件。能力题专门做一张逻辑连接词卡：不能辨认\emph{或者}不能控制。}

\section{第十一讲：犯罪形态}
\subsection{时间节点与停止原因共同决定评价}
令$t_p$为准备行为开始，$t_s$为着手，$t_f$为法定既遂条件实现。阶段划分不能只按行为人做了多少动作计算；着手与既遂均须依具体罪名解释。

\lawinput{刑法第22至24条分别规定犯罪预备、未遂与中止。未遂要求已经着手，因意志以外原因而未得逞；中止包括自动放弃犯罪或者自动有效防止犯罪结果发生。\lawref{L1}{第22至24条}}

令$S$为着手，$F$为达到该罪既遂，$V$为停止源于自动性，$E$为在需要防止结果的情形下已经有效防止。通常的分析树如下。
\begin{center}
\begin{tikzpicture}[>=Stealth,node distance=10mm and 11mm,every node/.style={draw,rounded corners,align=center,font=\small,inner sep=4pt}]
\node (q) at (0,0) {法定既遂条件是否实现？};
\node (done) at (-3,-1.4) {已经实现\\进入既遂评价};
\node (s) at (3,-1.4) {尚未实现\\是否已经着手？};
\node (prep) at (0,-2.9) {未着手\\预备或中止等评价};
\node (stop) at (5.1,-2.9) {已着手\\检查停止原因和有效性};
\draw[->](q)--node[above,draw=none]{是}(done);
\draw[->](q)--node[above,draw=none]{否}(s);
\draw[->](s)--(prep);\draw[->](s)--(stop);
\end{tikzpicture}
\end{center}

原书概览把“预备、中止、未遂、既遂”依次列示；这些词不能被理解为每个案件必经的状态转移链。中止与未遂等是需要根据停止时点、原因和效果区分的评价分支，中止还可能发生在预备阶段。

在着手后未既遂的案件中，若系意志以外原因而失败，进入未遂；若具有自动性并满足相应有效性要求，进入中止。只写$\neg F$不足以区分两者。因被害人反抗、工具失效、警察到场等停止是否属于意志以外原因，需要结合行为人认识和实际过程判断，不能只看事件名称。

\caseinput{甲对乙投毒后后悔，立即送医并如实告知毒物，医生据此救活乙。}
第一步，确认已着手；第二步，死亡未发生；第三步，确认甲的救助出于自动性；第四步，确认救助有效阻止死亡且与甲的行动有关；第五步，按中止规则评价，同时核查已经造成的其他损害。变式：甲虽后悔但拒不救助，乙靠无关人员偶然获救，不能直接把心理后悔等同于自动有效防止结果。

\subsection{处罚后果是不同的规范模态}
预备犯可以比照既遂犯从轻、减轻或者免除；未遂犯可以比照既遂犯从轻或者减轻；中止犯没有造成损害的，应当免除，造成损害的，应当减轻。以上需置于条文语境，不能将其统一成“未完成就减刑”。“免除处罚”与“不构成犯罪”也不是同一结论。

\memorycard{一条时间线加两个问题：“是否着手/既遂”“为什么停止/是否有效”。处罚卡单独背模态词，不把“可以”改成“应当”。}

\section{第十二讲：共同犯罪}
\subsection{共同事件与个人责任分层}
\lawinput{刑法第25条规定共同犯罪为二人以上共同故意犯罪，共同过失不以共同犯罪论处；主犯、从犯、胁从犯、教唆犯分别有相应条文。\lawref{L1}{第25至29条}}

以参与者为顶点建立有向图$G=(V,E)$，边表示教唆、帮助、共同实行等具体联系。图用于记录贡献，不能把“和某人认识”或“同处现场”画成犯罪联系。对每个人$p_i$，至少保留
\[
 (\text{参与行为}_i,\text{共同故意范围}_i,
 \text{介入时间}_i,\text{身份}_i,\text{责任条件}_i).
\]
在共同范围内研究行为贡献及归责，再对各参与者分别判断责任。共同故意不能从甲的口供单向复制到全部在场者。

\caseinput{甲乙共同计划盗窃，乙负责望风；甲入内后因私人仇恨另行杀害丙，乙既不知情也未提供针对杀人的帮助。}
第一步，确定盗窃的共同计划与乙的帮助。第二步，把杀人行为单独列出。第三步，检查杀人是否落在乙的共同故意与可归责范围内。第四步，不能仅因共同盗窃而当然使乙对甲的另起杀意负责。是否存在其他责任，应依额外事实另行判断。

\subsection{正犯、教唆、帮助与作用大小是两种分类轴}
参与方式回答“如何参与”，主犯从犯回答“在共同犯罪中起何种作用”。两个轴有联系，但不是简单同义词。不能机械规定“凡望风都是从犯”或“凡教唆都是主犯”。应以个案中的支配、组织、策划、实行和实际贡献等事实结合规范作判断。

对教唆，记$T(i,j,k)$为$i$使$j$形成实施$k$的犯罪决意。若$j$原已有确定决意，某句话是否仍属教唆需具体分析。对帮助，记$H(i,j,a)$为$i$对$j$实施犯罪提供具有相应意义的促进。主观明知的对象、帮助时点与实际作用均不可省略。

\subsection{承继、片面与身份问题}
后加入者介入时间$t_i$是关键。先列$t<t_i$已完成的事实与$t\ge t_i$继续实施的事实，再研究后者是否承接了仍在持续的危险或结果进程；不能把全部前史不加区分地转嫁。片面共犯、共犯从属性以及身份犯问题存在理论差异时，应把采用的观点写成$I$参数，并分别演算。

\begin{modelprop}[共同故意范围的集合检查]
设$S_i$是有证据支持的参与者$i$的共同故意范围。若某结果$r\notin S_i$，仅凭“$i$参加了同一事件”这一事实，不能用本模型推出其对$r$承担故意共犯责任。
\end{modelprop}
\begin{proof}
模型的故意共犯归责以$r\in S_i$为必要前提。已知$r\notin S_i$时，该前提不成立；“同一事件”与$r\in S_i$不是同义命题，不能替代。是否成立过失或其他责任是另一待检规则。
\end{proof}

\memorycard{每个参与者各写一行，列“做了什么、知道什么、想共同实现什么、何时加入”。禁止只画一张组织关系图便给所有人同一个结论。}

\section{第十三讲：罪数}
\subsection{先确定行为单元，再处理评价关系}
令$A=\{a_1,\ldots,a_m\}$为待评价行为单元，$K=\{k_1,\ldots,k_n\}$为候选罪名。矩阵
\[
 M_{ij}=1\quad\text{表示行为 }a_i\text{在给定前提下符合 }k_j\text{的候选构成}.
\]
列出$M$只是找出可能重叠的评价，不能直接把所有$1$相加当作最终罪数。接下来必须检查法条竞合、想象竞合、法律拟制、吸收或其他特定处理规则，以及应当数罪并罚的独立行为。

\caseinput{甲为杀害乙而开一枪，弹丸穿过乙又击伤丙。}
先确定行为单元与结果分别涉及何种故意或过失，再列候选构成；之后讨论一个行为触犯数罪名的处理。不能先把“一枪”当作“必然只成立一个罪名”，也不能按受害人数直接相加成罪数。

\subsection{特别规则优先，但必须有适用依据}
某些条文规定转化犯或按照另一罪定罪处罚。用映射$T_N:k\mapsto k'$记录法定转化时，应同时记录触发条件$P_N$：
\[
 C_k\wedge P_N\Rightarrow\operatorname{EvaluateAs}(k').
\]
没有$P_N$，不能只凭行为“相似”完成映射。转化型抢劫就是需要逐项审查前行为、目的、当场性及暴力胁迫等条件的例子。

对多个行为，先看是否相互独立，再查有没有法律上只作一罪评价的依据。评价同一事实时还要避免在构成、加重和量刑中无依据地重复使用同一不利因素；但也不能把“一次评价”误解成只要同一事件就只能适用一个规范。

\modelbound{本节矩阵是查漏工具，不是自动决定罪数的算法。行为单复数本身可能涉及规范评价，必须说明划分理由。}
\memorycard{固定顺序“行为单元—候选罪名—竞合关系—处理依据—最终结论”。每个择一重处结论旁写出具体依据，不能只背一个口号。}

\chapter{总论：刑罚、裁量、执行与消灭}

\section{第十四讲：刑罚的体系}
\subsection{刑罚是带类型的对象}
把一项刑罚写为$S=(\tau,d,a)$，$\tau$为刑罚种类，$d$为依法可用的期限、金额等参数，$a$为附加刑或附随处理。不同类型之间不能随意做加法；两年有期徒刑加一万元罚金，不能写成某个单一的“总惩罚数”。

\lawinput{刑法第33条列举主刑，第34条列举附加刑；附加刑也可以独立适用。具体适用须再查各刑种的条件和有关分则规定。\lawref{L1}{第33至60条}}

教学表可列为：主刑中的管制、拘役、有期徒刑、无期徒刑、死刑；附加刑中的罚金、剥夺政治权利、没收财产。对犯罪的外国人可适用驱逐出境的规定应另列，不能混入五种主刑之一。

\caseinput{某罪依法判处有期徒刑并处罚金。}
第一步，识别“并处”要求主刑与罚金的组合；第二步，分别核对主刑区间和罚金规则；第三步，确认该条是否允许单处罚金；第四步，不能因为缴纳罚金就把应执行的自由刑自动抵销。

\subsection{非刑罚措施不能因名称相近而混淆}
责令赔偿、训诫等非刑罚处理与刑罚应分别建模。财产返还、追缴、退赔与罚金的目的和依据也不同。把涉案财物全部相加称为罚金，会同时混淆计算对象与法律性质。

\modelbound{刑罚轻重可以有规范上的排序，但不是一切维度均能量化为同一数轴。死刑、无期徒刑与有期徒刑之间尤其不能用任意常数换算。}
\memorycard{先背类型，再背期限与适用条件。做题把“追缴、退赔、罚金、没收财产”分成四格，并说明每一格处理的是什么。}

\section{第十五讲：刑罚的裁量}
\subsection{从重、从轻、减轻与免除的区间差别}
设已经正确确定的法定刑幅度为$I=[l,u]$。从重与从轻原则上都在法定幅度内依法选择；减轻处罚涉及低于法定刑处罚，并须符合第63条等规则；免除处罚并不等于从来没有构成犯罪。\lawref{L1}{第61至63条}

学习模型只能给出可行集合$\mathcal S(N,F)$，不能仅凭几个情节加减得到唯一刑期。若设一条固定公式$y=\text{基数}-2\times\text{自首}+1\times\text{累犯}$，除非有对应的适用依据和范围，否则数字只是自造，不能作为法律推导。

\subsection{自首、坦白、立功和累犯各有独立条件}
自首应区分自动投案并如实供述的通常结构与法律规定的准自首等情形；坦白和立功也不能互相替换。\lawref{L1}{第65至68条} 一般累犯需要对前后犯罪性质、刑罚条件、时间、年龄等逐项核查，不能把“曾经犯罪”当作完整条件。对特别累犯另依规定。

\caseinput{甲如实供述了司法机关尚未掌握的本人另一罪行，同时揭发乙的犯罪。}
第一步，区分本人罪行与他人犯罪；第二步，查本人供述是否满足相应自首规则；第三步，对乙的揭发查是否属实并符合立功条件；第四步，分别确定规范效果，不能因甲“说了很多”就把两项条件混成一个从宽情节。

\subsection{数罪并罚的区间演算}
\lawinput{判决宣告前一人犯数罪，对于有期徒刑的通常并罚，执行刑期受数刑最高刑期、总和刑期和法定上限共同约束；总和不足35年与达到35年的上限不同。异种刑罚与附加刑另依规定。\lawref{L6}{第69条}}

只在数刑均为有期徒刑、无另行特别限制的本例范围内，令
\[
 m=\max_i s_i,\quad S=\sum_i s_i,\quad
 U(S)=\begin{cases}20,&S<35,\\25,&S\ge35.\end{cases}
\]
则执行刑期$y$需要满足
\[
 m\le y\le\min\{S,U(S)\}.
\]
该区间不是唯一量刑结果；其中的酌定仍须说明依据。法条中“以上”“以下”通常包括本数，应与刑法第99条等规定结合理解。

\caseinput{三罪分别判处8年、10年、12年有期徒刑。}
先算$m=12$，再算$S=30<35$，故$U=20$，得到$12\le y\le20$。若另例为10年、12年、15年，则$S=37$，上限为25年，故$15\le y\le25$。不能把30年或37年直接作为执行刑期，也不能声称区间内任一数都无须理由。

\subsection{漏罪与新罪不能交换运算顺序}
已执行刑期为$e$，原判刑为$s$，后判刑为$t$，并罚算子为$\mathcal C$。漏罪按法定规则体现为先将$s,t$并罚，再计入已执行$e$；新罪则以原刑未执行部分$s-e$与$t$并罚。\lawref{L6}{第70、71条}
\[
 \text{漏罪剩余}:\ \mathcal C(s,t)-e,
 \qquad
 \text{新罪之后}:\ \mathcal C(s-e,t).
\]
$\mathcal C$不是通常加法，两个表达式一般不能交换化简。缓刑还需独立审查适用刑罚、情节、悔罪表现、再犯危险、社区影响和排除条件，不能由刑期较短直接推出。

\memorycard{处罚模态四张卡：“范围内轻/重”“范围下减轻”“免除”“缓刑条件”。并罚题先写类型，再算最高、总和、上限，最后处理时间顺序。}

\section{第十六讲：刑罚的执行和消灭}
\subsection{减刑与假释是不同状态转换}
用状态记录$S_t=(\text{刑种},\text{原判},\text{已执行},\text{执行状态})$。减刑改变依法应执行的刑罚内容；假释是在法定条件与考验规则下提前释放，不能简单等同刑罚已经消灭。\lawref{L1}{第78至86条}

\caseinput{甲被假释后进入考验期。}
第一步，保留原刑与尚未执行部分的信息；第二步，记录考验期条件；第三步，依是否出现新罪、漏罪或其他法定撤销事由分别处理；第四步，不能在假释决定作出时就删除全部刑罚状态。对假释禁止条件与实际执行期限应单独核对。

\subsection{追诉时效的分段函数}
依刑法第87条，把法定最高刑$u$对应到追诉期限$T(u)$：
\[
 T(u)=\begin{cases}
 5\text{年},&u<5\text{年有期徒刑},\\
 10\text{年},&5\le u<10\text{年有期徒刑},\\
 15\text{年},&u\ge10\text{年有期徒刑},\\
 20\text{年},&\text{法定最高刑为无期徒刑或死刑}.
 \end{cases}
\]
这里第三档与第四档按刑种区分；20年以后认为必须追诉的，还需依法律规定核准。时效起算、继续或连续状态、期间内再犯罪，以及不受追诉期限限制的情形另依第88、89条分析。\lawref{L1}{第87至89条}

\caseinput{某罪适用的法定最高刑为7年有期徒刑。}
先定位$5\le7<10$，得到通常追诉期限10年；再确定起算日及是否发生法定特殊情形。只完成第一步，不等于已经算出本案最后追诉日。不能拿最终宣告刑2年替换法定最高刑7年作为分档依据。

\memorycard{执行卡用“原判—已执行—剩余—考验”四格。时效卡背“法定最高刑”，再写“起算、重算、不受限制、核准”，防止只会背5、10、15、20。}

\chapter{分论：罪名结构、人身与财产}

\section{第十七讲：分论概说}
\subsection{罪状是结构化规则，罪名是索引}
一个分则条文可拆为
\[
 N_k=(\text{主体},\text{行为},\text{对象},\text{结果或危险},
 \text{罪过},\text{入罪条件},\text{处罚规则}).
\]
罪名只是查找该结构的索引。看到名称相近，并不能把一个罪的条件借给另一个罪。空白罪状还需要查明被援引规范的内容、效力与适用时间；援引行政规范不意味着行政违法自动等于刑事犯罪。

\subsection{注意规定与法律拟制的区别}
注意规定用于提示本可由一般规则得出的处理；法律拟制则依法将某类事实按另一种构成评价。两者不能只凭文字是否出现“以某罪论处”机械区分，需要比较一般规则与特别条文的关系。

\caseinput{先建立一般规则$P\Rightarrow Q$。特别条文又写明某种$P_1\subseteq P$适用$Q$。另有条文将通常不满足$P$的$R$也依法评价为$Q$。}
前一种情形在该模型中没有扩张一般规则的对象；后一种需要特别依据$R\Rightarrow Q$。不能在没有特别条文时自行添加后一条箭头，也不能仅依据这两个集合关系就断言现实中某条一定属于哪一类。

\modelbound{分则适用还须核对总则的犯罪形态、共犯与阻却事由；孤立罪状不是完整定罪规则。}
\memorycard{每读一个罪名，用七格表填满条件。遇到“依照前款”“以某罪论处”“违反国家规定”，用箭头指向被援引的具体条文。}

\section{第十八讲：人身犯罪}
\subsection{生命、身体、性权利、自由与人格分别检索}
原书本讲涵盖生命身体、性权利、人身自由、名誉人格、婚姻家庭等主题。形式化不能把它们都缩成“伤害他人”。分别建立法益轴$L$和行为方式轴$A$，再匹配分则构成。

对于致人死亡，最先分开杀人故意、伤害故意、其他故意及过失；对于身体侵害，还需核对伤情、具体行为、结果归责和特别规定。\lawref{L1}{第232至235条} 对性权利犯罪，须按具体条文检查对象年龄、行为方式、强制或违背意志等条件及特别职责关系，不以道德评价代替要件。刑法修正案（十一）增加的负有照护职责人员性侵罪，有独立主体、对象与行为条件。\lawref{L2}{第27条}

\caseinput{甲意图使乙暂时失去行动能力，持续捆绑并殴打乙，乙死亡。}
第一步，把限制自由与伤害行为分别记录。第二步，查明甲对死亡和伤害的认识、态度。第三步，查明死亡机制及特别条文对暴力致伤、致死的处理。第四步，才决定故意伤害、故意杀人、非法拘禁有关结果加重或其他竞合评价。不能从“最初只想关住”就免除对后续行为的评价。

\subsection{非法拘禁与绑架：控制人的目的并非唯一差别}
两类案件都可能具有控制人身自由，但绑架需审查勒索财物目的或作为人质等法定结构，以及具体行为。为索取债务而拘禁他人，须核对法条特别规定与债务事实；不是所有索要金钱的控制行为都可以只按“要钱”二字分类。\lawref{L1}{第238、239条}

对侮辱、诽谤，应分别检索行为方式、对象、情节和告诉处理规则；对婚姻家庭犯罪，应核对特定身份关系、持续性、作为义务、严重程度等。不能把所有家庭冲突自动解释为刑事犯罪，也不能因存在家庭关系就排除犯罪可能。

\subsection{一个最小变式组}
同样是乙受伤，设三例：甲故意打伤；甲在符合条件的正当防卫中使乙受伤；甲完全无法预见且不可避免地造成乙受伤。结果变量相同，行为性质、主观要件与阻却条件不同，法律评价可能分别不同。这证明“伤情严重程度”不是足以覆盖全部定罪判断的唯一输入。

\memorycard{先背五组法益，再对每组背“行为—对象—主观—结果—特别规则”。生命身体题必须先写杀人故意与伤害故意的区分依据。}

\section{第十九讲：财产犯罪}
\subsection{财产题的五个基础变量}
用五元组描述取财过程：
\[
 (O_0,M,D,O_1,U).
\]
$O_0,O_1$分别为行为前的占有和行为后的控制状态；$M$为取财机制，$D$为处分及相关认识，$U$为非法占有目的。另列所有权、处分权限、价值和财产性利益，不能混称“钱是谁的”。

\begin{center}
\begingroup\small
\begin{tabular}{p{.15\textwidth}p{.35\textwidth}p{.39\textwidth}}
\toprule
候选类型 & 首要检索问题 & 仍须补齐\\\midrule
侵占 & 是否先有代为保管等法定关系及相应占有？ & 对象类型、非法占有、拒不退还或交出等法定条件\\
盗窃 & 是否违背占有人意愿转移占有？ & 非法占有目的、入罪路径、阶段与责任条件\\
抢夺 & 是否以夺取方式针对财物施力？ & 对人危险、特别规定和转化情形\\
抢劫 & 是否以暴力胁迫等法定方法取财？ & 方法作用、财物取得、时空与特别条件\\
敲诈勒索 & 是否通过威胁迫使他人交付财物？ & 非法占有目的、胁迫内容、入罪条件\\
诈骗 & 是否因欺骗导致错误认识及处分？ & 处分与损失的联系、主观目的及其他条件\\\bottomrule
\end{tabular}
\endgroup
\end{center}
本表只列检索入口。占有、处分意识、财产性利益等争议须标明解释立场，完整构成还须补齐。

\subsection{诈骗的因果链及断点}
在本册采用的通常教学框架下，记欺骗$D$、错误认识$E$、处分$P$、财产转移$T$、损失$L$，用
\[
 D\longrightarrow E\longrightarrow P\longrightarrow T\longrightarrow L
\]
表示待证明的机制链。箭头不是单纯时间顺序。每一段都要说明前项为何导致后项，另核对非法占有目的等条件。

\caseinput{甲假称出售真实存在且愿意交付的商品，乙因此付款；甲实际从开始就没有交货意图。}
第一步，核对甲关于交易能力、履约意愿等所作虚假表示。第二步，确认乙因该表示形成错误认识。第三步，确认乙基于错误认识付款。第四步，确认财产损失及非法占有目的。若只是履约后发生真实商业困难，第一步与主观条件可能改变，不能把所有违约都认定诈骗。

反例：乙已经识破骗局，仍为协助侦查而交付款项。错误认识环节并未按设例成立，不能照搬通常诈骗既遂的完整链；应另查未遂与其他事实。三角诈骗还须检查处分人是否具有处分他人财产的地位或权限及相应解释条件。

\subsection{盗窃的入罪条件是析取，金额不是唯一通道}
\lawinput{刑法第264条除数额条件外，还规定多次盗窃、入户盗窃、携带凶器盗窃、扒窃等路径。2013年司法解释细化了相关概念，并给出各地数额标准的确定框架。\lawref{L1}{第264条}\lawref{L7}{第1至3条}}

令$A$为数额达到适用标准，$M,H,W,P$分别为多次、入户、携带凶器、扒窃的法定情形，则入罪路径检索式为
\[
 A\vee M\vee H\vee W\vee P.
\]
该析取还须与盗窃行为、非法占有目的、责任等条件相结合。若$A=0$，只能排除数额这一通道，不能推出整个析取为$0$。在两年内三次以上的多次盗窃规则中，“次”的划分仍需按事实和规范判断，不能把同一过程切割成任意次数。

金额标准记为$\theta(\text{地区},t,N)$。对普通数额型案例，$v\ge\theta$是一个必要的金额检验，但不得用某地一个数字覆盖全国全部案件。对盗窃未遂，需再检索司法解释关于目标数额、珍贵文物和其他严重情节等规定，不能机械认为凡未取到就绝对不处罚。

\subsection{侵占、职务侵占与盗窃：先画占有变化}
侵占题重点检查先前依法或事实上形成的特定占有关系及后续非法据为己有，盗窃题关注转移他人占有；职务侵占再增加单位人员、职务便利、单位财物等条件。\lawref{L1}{第270、271条} 职务上的保管控制、工作场所的接近机会和一般员工身份不是同一个条件。

\caseinput{甲受托保管乙的相机，随后以所有人身份出售并拒绝返还；另例中，甲从乙仍控制的柜中拿走相机。}
第一例先检查代为保管及侵占的其他条件；第二例先检查是否打破乙的占有。若甲是公司员工，还须说明其接近财物是职务管理控制，还是仅有便利机会，不能仅凭工牌直接选职务侵占罪。

\subsection{抢劫、抢夺、敲诈勒索与转化}
对暴力胁迫，记录对象、强度、即时性、时空联系和取财机制。“感到害怕”不足以唯一区分抢劫与敲诈勒索。对携带凶器抢夺、盗窃诈骗抢夺后当场为特定目的使用暴力胁迫等情形，应逐项核查法定转化条件。\lawref{L1}{第263、267、269、274条}

\caseinput{甲实施盗窃后，为抗拒当场抓捕而使用暴力。}
先确认前行为符合转化规则所要求的范围，再确认目的属于窝藏赃物、抗拒抓捕或毁灭罪证等法定范围，继而判断当场性、暴力胁迫及相关条件，最后才讨论按抢劫评价。若数日后单纯因私人纠纷施暴，当场性与关联目的改变，不能自动套同一转化规则。

\subsection{既遂、金额、利益与罪数分开计算}
既遂应依罪名与对象审查占有、取得或损失，不能统一以“钱花出去”为标准。犯罪数额、被害人损失、违法所得、追缴额、赔偿额可能不同，分别记为$v_c,v_l,v_g,v_r,v_d$，避免一数多用。

财产性利益问题必须说明利益的内容、可评价性、处分或控制方式及具体罪名是否容纳。罪数问题则在各行为构成查明后处理：取财前的欺骗、取财中的暴力、事后的毁坏和掩饰行为，既不能一律并罚，也不能一律吸收。

\memorycard{财产题固定画三个时点的占有图，再问“怎样取得”“为什么交付”“何时形成非法占有目的”。六类罪名两两对比，每对只改变一个关键事实。}

\chapter{分论：公共利益、社会管理与职务犯罪}

\section{第二十讲：危害公共安全犯罪}
\subsection{公共危险必须说明开放范围}
危害公共安全不只是受害人数多。分析时记录危险所及人员或财产范围、扩散可能、可控制性与具体场景。针对特定个人的侵害，可能同时危及公共安全，也可能没有达到有关公共危险；需要根据方法与环境判断。

\lawinput{放火、决水、爆炸、投放危险物质等条文与相应过失犯罪、交通犯罪、安全生产犯罪各有不同构成。不得把所有危险行为直接装入一个兜底罪名。\lawref{L1}{第114至139条相关条款}}

\caseinput{甲为报复乙，在人员密集且有蔓延条件的场所点燃物品；变式为在完全隔离且无扩散可能的特定环境中损毁乙的物品。}
先查火势或危险的客观范围，再查甲的认识与态度，随后比较放火等公共安全犯罪与针对个人、财产的其他构成。两例均有“点火”，但公共危险这一条件可能不同。不能把动作名称直接当罪名。

\subsection{六类主题的统一检查表}
危险方法型重点查方法与公共危险；破坏型查设施对象、运行状态和危险；交通型查具体交通行为、责任、结果或危险标准；恐怖活动型查法律定义及具体参与行为；枪支类查物品性质、行为类型和数量情节；安全生产类查职责、违章或危险行为、结果与归责。每类的专业鉴定或行政标准都应作为明确输入，不能凭肉眼印象替代。

\modelbound{“其他危险方法”不能被无限扩展；是否具有与列举方法相当的危险性等问题，需要严格解释。安全生产中的一般管理过失与特定严重行为也不能混同。}
\memorycard{先写“危险波及谁、怎样扩散、能否控制”，再进入六类犯罪各自的对象与行为条件。}

\section{第二十一讲：破坏经济秩序犯罪}
\subsection{行政违法只是前置检索，不自动完成入罪}
经济犯罪经常涉及许可、资质、交易规则和监管秩序。设$A$为违反行政规范，$C$为刑事构成，通常不能推出$A\Rightarrow C$。刑事构成还可能要求特定主体、行为、数额、情节、危险或结果以及罪过。

本讲八组主题可用不同变量表达：伪劣商品的产品属性与销售额等；走私的监管边界与应税额等；公司企业管理秩序的职务与职责；金融管理的资金来源、募集方式与许可；金融诈骗的非法占有目的和欺骗机制；税收的税种、行为与法定处理条件；知识产权的权利有效性、许可范围及侵权行为；市场秩序的具体禁止行为和严重性。统一的是审查程序，不是所有罪名的门槛。

\caseinput{企业未经某项许可开展经营并造成交易损失。}
第一步，确认许可规范的层级与效力；第二步，确定违反的具体义务；第三步，查刑法是否将该类行为及情节规定为犯罪；第四步，查明主观条件和入罪门槛。只有第一步和第二步成立时，仍不能跳到刑事有罪结论。

\subsection{金额变量必须标注口径}
设销售金额$S$、违法所得$G$、偷逃税额$T$、给他人损失$L$。一般有$S\ne G\ne T\ne L$，不同条文使用的口径不能交换。若某条文要求$T\ge\theta$，把$S$代入即属于变量类型错误，即使两者单位都是“元”。

对空白罪状，应给援引规范添加版本。新司法解释改变数额认定或行为范围时，先处理时间效力，再更新参数；不能在旧案中无条件代入新阈值。刑法修正案（十二）涉及公司企业人员有关犯罪及行贿等条款的修订，须纳入版本检查。\lawref{L3}{第1至7条}

\modelbound{本节没有列全国各罪的全部最新金额表，因此不能把本册当成经济犯罪门槛的完整数据库。具体金额题须查当时适用的专项解释。}
\memorycard{每个数字旁必须写口径：“销售额、所得、税额、损失”。再添一个来源和生效时间，避免只背裸数字。}

\section{第二十二讲：妨害社会管理犯罪}
\subsection{从九组主题提取各自关键条件}
扰乱公共秩序、妨害司法、文物、国边境、公共卫生、环境、毒品、卖淫及淫秽物品九组主题，共享“社会管理”这一章名，但不共享一套充分必要要件。学习时必须恢复各自对象与行为。

妨害司法题记录所处诉讼阶段、人员身份、证据或执行对象和具体妨害方式；文物题记录法律属性、等级、来源与行为；国边境题记录出入境行为及组织运送关系；公共卫生题记录法定管理措施、对象和传播风险；环境题记录污染物、处置方式、因果与严重程度；毒品题记录物质鉴定、数量、行为类型与明知；卖淫、淫秽物品题分别核对组织、强迫、介绍、传播等具体行为及对象。\lawref{L1}{分则第六章相关条款}

\caseinput{甲替乙保管一个包裹，包内检出毒品。}
第一步，物质鉴定支持对象性质；第二步，保管事实支持客观接触；第三步，仍需查甲是否明知及具体参与行为；第四步，不能用“检出毒品”补足主观条件。若另有充分证据证明明知，还需区分具体行为类型，不能把所有接触都叫贩卖。

\subsection{不能把计量便利当成法律定义}
数量是常见分档变量，但具体犯罪是否处罚、如何认定既遂、是否需要数量门槛，取决于罪名。对于明知的证明，也不能凭一个任意概率阈值，例如“有70\%可能知道”就认定刑事故意。证据法要求与统计模型中的数值信念不是同一对象。

环境与公共卫生案件还常有多源作用和专业因果问题。把检测值超过某一行政标准与已经满足刑法所有要件等同，可能遗漏行为、主体、严重程度和归责条件。应按“规范—检测—行为—后果—主观—责任”依次审查。

\memorycard{九组各做一张“对象—动作—门槛—主观”卡。毒品、文物、污染物等对象性质单独列证据来源，防止把物理接触当成全部构成。}

\section{第二十三讲：贪污贿赂犯罪}
\subsection{职务犯罪至少有三条独立的线}
第一条是身份线：行为人是否属于相应主体。第二条是职务线：行为与职务便利、职权地位或影响力如何连接。第三条是财产与利益线：财物从谁流向谁、为谁谋取何种利益、归谁所有。任何一条缺失都不能用其他两条自动补齐。

贪污与挪用公款应区分非法占有与挪用用途、数额、时间等条件；受贿与行贿应分别判断，不以一方最终定罪作为另一方全部条件已经证明的替代；斡旋受贿、利用影响力受贿和介绍贿赂则要精确记录关系与行为方式。\lawref{L1}{第382至396条}\lawref{L3}{有关行贿条款修订}

\caseinput{甲为国家工作人员，收受乙财物，并承诺通过丙的职务行为为乙谋取利益。}
第一步，画出乙向甲的财物流向；第二步，说明甲自身职权地位与丙之间的联系；第三步，查谋取的利益以及承诺的内容；第四步，依具体关系检索直接利用职务便利、斡旋受贿或其他规则。不能只因有第三人丙就一律改认介绍贿赂。

\subsection{2026年解释的版本化输入}
\lawinput{最高法、最高检于2026年发布贪污贿赂司法解释（二），自2026年5月1日起施行。该解释涉及单位犯罪标准、有关民营企业人员犯罪的标准参照、挪用公款、预期收益型受贿、职务联系与介绍贿赂等问题。\lawref{L8}{施行时间及第1至19条}}

以下只以单位受贿的一个分档规则演示形式化。设$v$为依适用规范认定的数额，$E$为该解释第1条列举的至少一项特别情形，单位受贿“情节严重”的相应金额条件为
\[
 Q(v,E)=(v\ge20\text{万元})
 \vee(10\text{万元}\le v<20\text{万元}\wedge E).
\]
这只是条文中的情节条件，不是单位受贿罪的完整定义，主体、行为、主观等仍需另查。\lawref{L8}{第1条}

\caseinput{其他构成条件均已明确，某单位受贿数额为15万元，并存在第1条所列的一项特别情形。}
第一步，$15<20$，第一条金额路径不成立；第二步，$10\le15<20$且$E=1$，第二路径成立；第三步，析取$Q=1$。变式：$E=0$时，$Q=0$；不能以“接近20万元”补足该项门槛。此处不据$Q=0$概括排除其他法律规定的所有可能评价。

\subsection{受贿数额不是任意时间的最高估值}
财物收受、预期收益、案发实际获利、市场价与支付价等，应按解释规定确定口径。用于练习的差额式$V_m-V_p$，只有在对应的规范情形成立时才可使用；不能拿所有受贿财物的历史最高价减去零作为统一标准。\lawref{L8}{第11、12条}

2026年解释第8条规定有关非国家工作人员受贿、行贿、职务侵占、挪用资金标准的参照关系；这提醒我们，旧教材中单独背下的数字可能需要更新。参照关系是规范输入，不是因为“两个罪看起来差不多”而自行建立的数学同构。

\modelbound{本节完成了若干核心规则的结构化和2026年新解释的定位，并未逐条复核原书本讲所有历史案例及全部数额。原文观点与新规范出现差异时应逐案记录，不隐去差异。}
\memorycard{画三线：“身份”“职务关系”“财产利益流向”。每遇金额先查规范版本，再算金额，最后核对金额之外的构成条件。}

\section{第二十四讲：渎职犯罪}
\subsection{职责关系与归责是关键连接}
以$D$表示具体职责，$V$表示违反职责的行为，$R$表示法定损害或危险，$C$表示职责违反与结果的归责联系，$G$表示该罪要求的主观要件。核心检索模板为
\[
 \operatorname{Status}\wedge D\wedge V\wedge R\wedge C\wedge G.
\]
某人是机关工作人员、某单位发生重大损失，两项事实之间仍缺少职责与行为连接。不能用职位高低自动补足归责。

\lawinput{滥用职权、玩忽职守及其他特别渎职犯罪应依各自主体、行为、结果和罪过要求判断；有特别规定的，还需处理与一般条文的关系。\lawref{L1}{第397至419条相关规定}}

\caseinput{负责特定安全检查的甲，在有能力履行的情况下完全不检查，后发生严重事故；同时乙虽在同一机关工作，但职责不包括该项检查。}
第一步，确定甲乙各自职责，而不是仅看办公地点；第二步，列出具体不履行行为；第三步，查正常履行能否发现并避免相应危险及实际事故机制；第四步，分别审查主观要件与法定结果；第五步，对甲乙分别作结论。不能将单位事故的总损失平均分配给全部人员以构造刑事责任。

\subsection{受贿与渎职可能同时出现}
财物收受与职务违反是不同的事实轴。即使查明受贿，也需单独证明渎职构成；即使没有财物往来，也不当然排除渎职。二者同现时，依有关规定处理罪数，不能未经检索就宣布一律吸收或一律并罚。

\memorycard{职责题闭卷写出“谁负何责—如何违反—造成什么—怎样归责—主观何在”。对同一事故至少比较一名有该职责者与一名没有该职责者。}

\section{第二十五讲：危害国家安全犯罪}
\subsection{抽象标签不能替代具体犯罪行为}
危害国家安全犯罪需要按具体条文核对行为、对象、主体、目的、联系及情节等。国家安全的重要性不能成为跳过罪刑法定和证明要求的理由。对涉及国家秘密或情报的犯罪，还须依据适用规范准确确定对象性质，不能以“敏感”二字代替法律认定。\lawref{L1}{分则第一章，第102至113条}

\caseinput{甲向境外人员发送一份资料，资料性质与甲的认识均有争议。}
第一步，确认资料是否符合有关条文的法定对象；第二步，确认接收方、发送方式与具体行为；第三步，查明主观认识与条文要求；第四步，核对其他构成及程序要求。仅凭“发送到境外”不能自动证明所有犯罪条件；反过来，仅凭资料容易获得也不能免除对其法律性质的核查。

\subsection{关系图应有证据来源}
若用关系图记录组织、人员与联络，节点和边都应标明所依据的事实。统计相关、共同联系人、同一平台出现，不能直接等同具体犯罪联络。对身份、指令、行为与主观认识，应分别列证据，避免把一张密集网络图误作已经完成的证明。

\modelbound{本节只规定严格审查的方法，不将宏观政治判断或未经核实的身份标签作为事实输入。具体案件应以有效法律及证据为基础。}
\memorycard{按“哪一条—什么对象—什么行为—什么主观—何种证据”逐项检索。任何一步只有形容词而没有事实时，标记为待补。}

% Supplement: explicit semantics, derivations, countermodels, and retrieval practice.
\chapter{严格推导补篇：证明、反模型与解释分歧}

\section{把“推出”写成可以检查的定义}
前面二十五讲提供法律主题的模型入口。本章补上“严格”二字需要的语义基础。以下定义用于核查推理，不能取代法条、证据规则或者解释论证。

固定一个具体问题，设有限原子命题集合为$\mathcal P=\{p_1,\ldots,p_n\}$。原子命题应足够具体，例如“甲在实施行为时已满16周岁”，不能只写“甲符合条件”。一个\textbf{赋值}$v$将每个原子映射为$0$或$1$，再按通常命题逻辑解释合取、析取、否定与蕴涵。

把法律前提$N$、已明确选择的解释方案$I$、已经给定的事实$F$合在一起，记为$\Gamma=N\cup I\cup F$。定义
\[
 \mathcal M(\Gamma)=\{v:\text{$v$使$\Gamma$中每个命题均为真}\}.
\]
这是该组前提仍容许的模型集合。由此定义
\[
 \Gamma\models Q
 \quad\Longleftrightarrow\quad
 \forall v\in\mathcal M(\Gamma),\ v(Q)=1.
\]
意思是：在每一种符合这些前提的赋值中，结论$Q$都成立。实际写论证时还应要求$\mathcal M(\Gamma)\ne\varnothing$，并注明每项事实的证据来源。若前提自相矛盾，不能利用空集合上的全称命题给任意定罪结论盖章。

\begin{modelprop}[反模型足以否定逻辑蕴涵]
若存在一个$v\in\mathcal M(\Gamma)$，使$v(Q)=0$，则$\Gamma\not\models Q$。
\end{modelprop}
\begin{proof}
由定义，$\Gamma\models Q$要求集合中每一个赋值均使$Q$为真。所给$v$属于该集合，却使$Q$为假，因此全称条件不成立。
\end{proof}

\caseinput{前提只有“若成立故意杀人既遂，则存在死亡结果”，另知“存在死亡结果”。拟推出“成立故意杀人既遂”。}
令$K$为故意杀人既遂，$D$为死亡结果。前提为$\{K\Rightarrow D,D\}$。取$v(K)=0,v(D)=1$，两项前提均为真，结论为假。故结论不由这些前提推出。过失致死、自然死亡等可以帮助理解缺口，但反模型本身已经足以证明此处的形式错误；具体哪种法律评价成立，仍须另行审查。

\memorycard{遇到“只要有结果，就一定有某罪”，将待检验主张写成$D\Rightarrow K$，与已知规则$K\Rightarrow D$分开。尝试$K=0,D=1$，说明为何它满足已知前提却否定拟推出的结论。}

\section{不完整信息：不能把问号偷偷改成零}
固定同一组规范$N$和解释$I$，设$F$只给定部分事实，记$\Gamma_F=N\cup I\cup F$。$\mathcal M(\Gamma_F)$中通常有多个赋值，可以把结论分类为
\[
 \operatorname{Status}_F(Q)=
 \begin{cases}
  1,&\forall v\in\mathcal M(\Gamma_F),\ v(Q)=1,\\
  0,&\forall v\in\mathcal M(\Gamma_F),\ v(Q)=0,\\
  ?,&\text{上述两种情况均不成立。}
 \end{cases}
\]
这里先要求$\mathcal M(\Gamma_F)\ne\varnothing$。这种“考察所有一致补全”的方法，能够保留同一个原子在不同位置出现时的相关性。

\begin{modelprop}[补充一致事实只会缩小模型集合]
若$F\subseteq F'$，则$\mathcal M(\Gamma_{F'})\subseteq\mathcal M(\Gamma_F)$。若两者均一致且$\Gamma_F\models Q$，则$\Gamma_{F'}\models Q$。
\end{modelprop}
\begin{proof}
任取$v\in\mathcal M(\Gamma_{F'})$。由于规范与解释保持不变，且$F$中每项也在$F'$中，$v$必满足$\Gamma_F$，故$v\in\mathcal M(\Gamma_F)$。若$Q$在较大的$\mathcal M(\Gamma_F)$中处处为真，在其子集上也处处为真。
\end{proof}

适用这条命题须固定同一组规则和解释。发现新证据后撤销了一个先前认定，属于\emph{修正前提}，不一定满足$F\subseteq F'$；新法改变规则，也不是单纯增加事实。因而不能据此宣称“案件知道得越多，原有有罪结论就越不会改变”。

\subsection{三值表的准确边界}
本册开头的三值合取表适合检查必要条件缺口，但不能不加限定地替代所有公式的经典语义。例如，尚不知$P$真假时，逐格代入强三值运算会把$P\vee\neg P$记为“?”；按一致补全观察，每个补全都使该式为真。两个方法回答的细节不同。含重复原子、互斥条件或关联约束时，宜直接考察完整模型，避免把相关的未知项当作可以任意独立变化。

\caseinput{甲是否自动放弃尚不明确，其他中止条件已给定。}
一组一致补全使自动性成立，另一组使它不成立，因此不能证明中止结论。这里既没有证明甲自动，也没有证明甲不自动。应回到停止原因、行为人的认识和外部条件，而不是用一个问号结束事实调查。

\section{必要、充分、等价：逐一提供证明义务}
设$C$为一个明确限定的法律结论，$P$为某项条件。
\begin{center}
\begin{tabular}{p{.18\textwidth}p{.28\textwidth}p{.43\textwidth}}
\toprule
声称 & 必须证明 & 一个否定它的反模型\\\midrule
$P$是必要条件 & $C\Rightarrow P$ & $C=1,P=0$\\
$P$是充分条件 & $P\Rightarrow C$ & $P=1,C=0$\\
$P$与$C$等价 & 两个方向分别成立 & 上述任一反模型\\
$P$只提供线索 & 不承诺充分或必要 & 仍须列出待补条件\\\bottomrule
\end{tabular}
\end{center}

若把一个规则写作$C\leftrightarrow P_1\wedge\cdots\wedge P_n$，就同时声称“列出的每一项均必要”以及“这些项合在一起已充分”。这比只写$C\Rightarrow P_j$强得多。许多刑法口诀只适合做检索线索，不能自动升级为双向等价式。

\begin{modelprop}[否定一个析取分支不等于否定整体]
设背景前提$\Gamma$一致，$Q=A\vee B_1\vee\cdots\vee B_m$，且$m\ge1$。若对某个$j$，$\Gamma\cup\{\neg A,B_j\}$仍一致，则$\Gamma\cup\{\neg A\}\not\models\neg Q$；若增加任一$B_j$为真这一前提，则可推出$Q$。
\end{modelprop}
\begin{proof}
由相容性前提，可取一个同时满足$\Gamma$、$\neg A$与$B_j$的赋值。它使$Q$为真，因而使$\neg Q$为假，是第一项推断的反模型。第二项依析取引入：任何一个分支为真，整个析取为真。这里不假定各分支任意独立；若背景规则已排除所有替代分支，就不能再使用这个反模型。
\end{proof}

盗窃罪的数额路径与多次、入户、携带凶器、扒窃等路径正适合用此命题检查。\lawref{L1}{第264条}\lawref{L7}{第1至3条} 证明的是路径表达式的逻辑性质；各分支是否实际满足及是否具备其他构成条件，不能省略。

\section{最小充分条件集：把背诵变成可拆解结构}
给定一致的有限条件集$\Gamma$与目标$Q$，若$S\subseteq\Gamma$满足$S\models Q$，称其为一个支持集。一致性保证这里没有利用矛盾前提的空真来制造支持。若对每个$s\in S$都有$S\setminus\{s\}\not\models Q$，则称$S$为\textbf{包含意义上的最小支持集}。这里的“最小”不表示所有支持集中元素最少，也不自动表示法律上不存在其他规则路径。

\begin{modelprop}[纯合取的逐项缺失反例]
设$p_1,\ldots,p_n$彼此独立，$C=p_1\wedge\cdots\wedge p_n$。事实集$S=\{p_1,\ldots,p_n\}$是$C$的最小支持集。
\end{modelprop}
\begin{proof}
首先，$S$使每个$p_i$为真，依合取引入得到$C$。其次，任删$p_j$，都可令其为假并令其余$p_i$为真；这个赋值满足剩余事实，却不满足$C$。故删除任一项后均不能推出$C$。
\end{proof}

独立性假设不能忽略：若规范或事实已蕴含$p_1\Rightarrow p_2$，则$p_2$可能不是这份证明记录中必须单独给出的事实，虽然它仍是法律构成的一项必要条件。学习时应区分\emph{法律要件}与\emph{证明该要件所需的事实材料}。

\memorycard{一张规则卡背面写全条件后，为每个条件各做一个“仅删此项”的变式。若删掉某项仍能沿另一规则路径得到同一结论，就记录那条路径，不强行宣称条件绝对必要。}

\section{争议学说作为参数：分别推导再比较}
解释争议不能通过偷偷选一个定义而消失。设两种方案分别为$I_1,I_2$，得到
\[
 \mathcal C_1=\{Q:N\cup I_1\cup F\models Q\},\qquad
 \mathcal C_2=\{Q:N\cup I_2\cup F\models Q\}.
\]
交集$\mathcal C_1\cap\mathcal C_2$表示两方案共同支持的结论；对称差
\[
 \mathcal C_1\mathbin{\triangle}\mathcal C_2
 =(\mathcal C_1\setminus\mathcal C_2)\cup(\mathcal C_2\setminus\mathcal C_1)
\]
定位方案导致的结论差异。此运算不能告诉我们哪种解释更有法律依据；选择方案仍须比较法条、体系、规范目的及相应学说论证。

\caseinput{一个取财案件中，其他条件固定，争点只在交付者交付财物时是否具有相应处分认识。}
先把两方案对处分认识的要求各写成一个谓词。再分别代入同一事实，不改变其他变量。若一个方案认可诈骗链中的处分环节，另一方案不认可，差异应定位到这一环节；不能把整道题说成“任何答案都可以”，也不能为使结论一致而临时修改事实。

\begin{modelprop}[共同结论对所列方案的稳定性]
若所列每一组前提$N\cup I_i\cup F$均一致并均推出$Q$，则$Q$在这组方案范围内稳定。仅凭这一有限比较，不能一般性地保证所有未列方案也支持$Q$。
\end{modelprop}
\begin{proof}
第一部分是有限全称命题的直接展开。为否定第二部分所述的一般性保证，取一个原子$Q$，令$N=F=\varnothing$，所列每个$I_i=\{Q\}$，另设$I_*=\{\neg Q\}$。各组前提均一致，已列方案都推出$Q$，而$I_*$不推出$Q$。这是一般推断的反例，并非声称对任意具体$Q$都能找到相反方案；若$N\cup F$本身已经推出$Q$，那么所有一致的补充解释都会保留该结论。
\end{proof}

\section{从推导树到可执行检查}
一份完整记录至少保留三类节点：\textbf{依据节点}写明条文或解释方案，\textbf{事实节点}写明证据或题目假定，\textbf{推论节点}写明所用规则与上游节点。下面画的是论证依赖关系，不是原书图的替代品。

\begingroup\setlength{\topsep}{3pt}
\begin{center}
\begin{tikzpicture}[>=Stealth,every node/.style={draw,align=center,font=\small,inner sep=5pt}]
\node (r) at (0,0) {规范前提\\$A\wedge B\Rightarrow C$};
\node (a) at (-3.5,-1.3) {事实节点\\$A$};
\node (b) at (3.5,-1.3) {事实节点\\$B$};
\node (ab) at (0,-1.3) {合取引入\\$A\wedge B$};
\node (c) at (0,-2.6) {肯定前件\\$C$};
\draw[->](a)--(ab);\draw[->](b)--(ab);
\draw[->](ab)--(c);\draw[->](r.east)--(5.1,0)--(5.1,-2.6)--(c.east);
\end{tikzpicture}
\end{center}
\endgroup

检查程序可以枚举小型布尔模型、核算金额边界、寻找反模型，或者核对每个推论是否有父节点。程序不能仅凭上述图确定$A$是否有证据支持，也不能证明引用的规则是当前适用的完整法律规范。本工程提供的模型检查脚本属于这种有限检查，不是全部刑法已经机器证明的声明。

\chapter{完整演算补篇：边界、归责与复习路径}

\section{正当防卫：先确定分支，再做真值运算}
本节对应第九讲。先固定行为具有防卫性质、不属于滥用防卫权等排除情形，再判断是否符合特殊防卫的全部法定条件。一般防卫过当要求“明显超过必要限度”和“造成重大损害”同时成立。\lawref{L1}{第20条}\lawref{L4}{限度与特殊防卫}

令$B$表示已通过前述基础审查，$S$表示符合特殊防卫条件；在一般防卫分支中令$X$为明显超过必要限度，$H$为造成重大损害。仅为组织检索，定义四个分支：
\begin{align*}
 D_0&=\neg B,\\
 D_1&=B\wedge S,\\
 D_2&=B\wedge\neg S\wedge X\wedge H,\\
 D_3&=B\wedge\neg S\wedge\neg(X\wedge H).
\end{align*}
$D_0$要求继续按其他规则评价；$D_1$进入特殊防卫审查后的结论；$D_2$满足所建模的一般防卫过当双条件；$D_3$表示这两个条件未同时具备。特别地，$D_3$不能被随意命名为“行为绝未超过限度”，因为它包括$X=1,H=0$的赋值。

\begin{modelprop}[防卫检索分支的互斥与穷尽]
在$B,S,X,H$均已取确定真假值的模型内，$D_0,D_1,D_2,D_3$恰有一个成立。
\end{modelprop}
\begin{proof}
若$B=0$，只有$D_0$成立。若$B=1$而$S=1$，只有$D_1$成立。若$B=1,S=0$，再按$X\wedge H$的真假分类，分别只有$D_2$或$D_3$成立。各对分支在$B$、$S$或$X\wedge H$至少一项上要求相反，故不会同时成立。
\end{proof}

\caseinput{题目已确认基础条件$B=1$、不属于特殊防卫$S=0$；重大损害$H=1$，但$X$没有证据结论。}
第一步，排除$D_0,D_1$。第二步，保留两个一致补全：$X=1$落入$D_2$，$X=0$落入$D_3$。第三步，因结论分支不同，不能据$H=1$完成防卫过当的推导。第四步，补查侵害的性质、强度、紧迫性、双方力量等与限度有关的事实，不能直接用损害严重程度替代$X$。

\memorycard{先背分支前提$B,S$，再背双条件$X,H$。闭卷画出四分支，并回答：为什么$X=1,H=0$不能被语言误写成“没有超过限度”？}

\section{犯罪形态：时间、终点和停止原因不能压成一个轴}
本节对应第十一讲。给定罪名$k$，将其特定既遂标准写为$F_k$，着手标准写为$S_k$，不得用一个统一的“财产到账”或“动作做完”替代所有罪名。\lawref{L1}{第22至24条}

本式只用于故意犯罪既遂前的犯罪过程中，行为人已放弃本次犯罪（非暂时停顿）的情形；既遂后的返还、救助不属本式适用域。令$V$为自动性，$R$为当时还需要有效防止结果，$E$为所要求的有效防止。在这个已限定的域内，自动中止路径的检索条件可写成
\[
 Z=V\wedge(\neg R\vee E).
\]
这一式子把“有的情形仅停止实行即可”“有的情形必须有效防止结果”分开。$R$的判断依犯罪进程与具体规范确定，不能由行为人自行宣布。

\begin{modelprop}[需要防止结果时，仅有后悔不充分]
在$R=1$的分支中，$Z$化简为$V\wedge E$。因此$V$本身不是$Z$的充分条件。
\end{modelprop}
\begin{proof}
代入$R=1$，得$\neg R\vee E=0\vee E=E$，故$Z=V\wedge E$。赋值$V=1,E=0$使$V$成立而$Z$不成立，即为充分性命题的反模型。
\end{proof}

在本式适用域内，也不能从$Z=0$直接推出未遂；还须确认已着手、未得逞、意志以外原因等相应条件。域尚未确定时，应先区分信息不足、行为尚未停止、已经达到法定既遂等情形，不能一概给它们赋值$Z=0$。对全体案件把“既遂、中止、未遂、预备”写成未经限定的四选一函数，会掩盖这些前提。

记忆时画两个互相配合的表：横轴表记录准备、着手、实行终了、既遂等实际时点，不假定所有罪名均按同一顺序出现；条件表记录自动性、停止原因、是否需要防止结果及有效性。每次只改变一个变量，观察评价变化，避免把时间更晚误当成处罚当然更重。

\section{共同犯罪：逐人逐项构造归责矩阵}
本节对应第十二讲。设参与者集合$P=\{p_1,\ldots,p_m\}$，待归责行为或结果集合$R=\{r_1,\ldots,r_n\}$。分别建立
\[
 A_{ij}=\text{存在可论证的行为贡献},\quad
 I_{ij}=\text{属于该人的相应共同故意范围},\quad
 C_{ij}=\text{符合其他归责与责任条件}.
\]
这里仅检查须由共同故意覆盖的具体行为或结果；若某项加重结果依法允许过失归责，就必须另设相应主观条件，不能原样套用这一矩阵。在此范围内，故意共犯归责候选为$G_{ij}=A_{ij}\wedge I_{ij}\wedge C_{ij}$。此矩阵把个人条件分开，不能把一行复制给所有人。共同过失不以共同犯罪论处等规则仍须独立适用。\lawref{L1}{第25至29条}

\begin{modelprop}[共同范围不同，归责行可以不同]
即使两个参与者对同一结果均有行为贡献，只要其中一人的相应故意必要条件为$0$，两人的故意共犯候选结论就不能由行为贡献相同而被强行视为相同。
\end{modelprop}
\begin{proof}
给出反模型：设$A_{1j}=A_{2j}=1$，$I_{1j}=1,I_{2j}=0$，并取$C_{1j}=1$。则$G_{1j}=1$；无论$C_{2j}$为何，$G_{2j}=1\wedge0\wedge C_{2j}=0$。故行为贡献列相同不保证归责列相同。这不表示两行必然不同：若$C_{1j}=0$，两项候选也可能同为$0$。
\end{proof}

\caseinput{甲乙共谋盗窃，乙望风；甲另起杀人故意实施杀人，乙对此不知情，也无针对杀人的帮助。}
建两列“盗窃”“杀人”，建两行“甲”“乙”。先填共同盗窃的参与和故意，再单列甲的超出行为。对乙的杀人列不能从“共谋盗窃”补出$I_{2,\text{杀人}}=1$。若题目另增加乙事前支持杀人或当场有意加入等事实，应重新审查该列，而非沿用旧输入。

\modelbound{矩阵没有解决片面共犯、承继共犯及身份犯等全部解释争议；这些问题仍需固定$I$后分别计算。矩阵的价值是迫使论证写明“谁对什么负责”，而不是自动确定每个人罪名。}

\section{数罪并罚：上下界可算，具体刑期不能只靠加法}
本节对应第十五讲。仅讨论已确定的若干有期徒刑在刑法第69条相应规则内的并罚区间；其他刑种、特别规定及法定情节另行处理。\lawref{L1}{第69至71条}
设各罪刑期为$t_1,\ldots,t_n>0$，令$S=\sum_i t_i$、$M=\max_i t_i$，上限函数为
\[
 U(S)=\begin{cases}20,&S<35,\\25,&S\ge35.\end{cases}
 \qquad
 \mathcal Y=\bigl[M,\min\{S,U(S)\}\bigr].
\]
$\mathcal Y$记录本模型给出的刑期数值上下界，不表示其中每一个实数都可以作为合法宣告刑；具体计算单位及其他适用要求仍须核对。实际执行刑期$y$仍需依法裁量，不能从该区间凭数学唯一确定一个数。

\begin{modelprop}[区间非空的准确条件]
对给定数值，$\mathcal Y\ne\varnothing$当且仅当$M\le\min\{S,U(S)\}$。由于$M\le S$，该条件等价于$M\le U(S)$。
\end{modelprop}
\begin{proof}
实数闭区间$[a,b]$非空当且仅当$a\le b$，代入即可。$M\le S$由所有$t_i>0$得到，因此还需检查的限制只有$M\le U(S)$。
\end{proof}

若一个输入让下界大于上界，应先检查是否把不属于这一简单模型的刑种、既有并罚刑或特别规则硬塞进来，不得把空区间“修正”为任意期望刑期。模型适用性优先于代数运算。

\caseinput{三个已判刑期分别为12、10、8年。}
$S=30,M=12,U(S)=20$，因此$12\le y\le20$。若改成15、12、8年，$S=35,M=15,U(S)=25$，因此$15\le y\le25$。跨越35年这个分界改变的是上限函数，不能把所有情形统一套成20年。

\subsection{先并后减与先减后并为什么不能交换}
设$B(a,b)$只表示在具体规范与裁量下确定的并罚结果。漏罪与新罪适用的法律次序不同；例如符号形式$B(a,b)-e$与$B(a-e,b)$中，$e$是已经执行部分。即使暂取一个简化数值算子$B(a,b)=\max(a,b)$用于展示顺序，$a=10,b=8,e=3$时两式分别为$7$和$8$。这只是“不应默认两种运算可交换”的反例，\textbf{不是实际并罚算法}；真正计算必须依第70、71条及案件条件。

\memorycard{刑期卡先写适用范围，再写$S,M,U(S)$，最后才算区间。次序卡各画一条箭头，并在箭头上写对应法条；不能只背“加减顺序”而不认案件是漏罪还是新罪。}

\section{财产犯罪：状态转换与取得机制分别建模}
本节对应第十九讲。设财物$o$在时点$t$的占有关系为$\operatorname{Pos}(p,o,t)$。占有是具有规范评价内容的关系，不等同于GPS坐标；所有权记为$\operatorname{Own}$，处分地位记为$\operatorname{Disp}$，三者不得混用。

一次取财分析写成
\[
 (s_0,a,m,s_1,u),
\]
其中$s_0,s_1$分别记录行为前后的权利与占有状态，$a$是行为，$m$是使状态改变的机制，$u$是相应主观条件。仅保留起点与终点会丢失$m$；盗窃、诈骗、抢劫等都可能出现财物从乙转入甲控制的外观，区别必须回到取得机制及完整构成。\lawref{L1}{第263、264、266、270、271条}

\begin{modelprop}[相同终点不能唯一识别取得机制]
设两个机制$m_1\ne m_2$均可把$s_0$变为$s_1$。仅由有序对$(s_0,s_1)$不能在两者之间唯一确定机制。
\end{modelprop}
\begin{proof}
若存在仅依赖$(s_0,s_1)$且总能正确识别机制的函数$f$，它面对两种过程都必须给出同一输出；正确性却分别要求输出$m_1$与$m_2$。因$m_1\ne m_2$而矛盾。
\end{proof}

\caseinput{过程甲：行为人平和地秘密拿走他人控制中的手机；过程乙：行为人欺骗占有人，使其陷入错误并处分同一手机。其余相关条件按各自构成另行补齐。}
两过程的最后控制状态可相同，但取得机制不同。因此“手机最后在谁手里”只能回答结果问题，不能完成盗窃与诈骗的区分。

\subsection{三角诈骗的关系表}
把欺骗者$d$、受骗人$p$、财产受损者$b$分开；本小节用$b$表示人，避免与前述财物$o$混淆。关系表至少有：欺骗$d\to p$、错误$p$、处分$p\to\text{财产}$、处分地位$p\rightsquigarrow b$、损失$b$。在采用要求相应处分地位的解释方案时，一项机制检查写为
\begin{align*}
 K_I(d,p,b)={}&D(d,p)\wedge E(p)\wedge P(p,b)\wedge A_I(p,b)\wedge L(b)\\
 &\wedge C_{DE}(d,p)\wedge C_{EP}(p,b)\wedge C_{PL}(p,b).
\end{align*}
其中$P(p,b)$表示$p$处分归属于受损者$b$的财产；$C_{DE}$、$C_{EP}$、$C_{PL}$分别要求有关欺骗引起有关错误、该错误引起该处分、该处分造成相应财产损失。仅把这些事件同时列为真，并不能替代因果与规范联系的论证。本式只检索这条机制，并非诈骗罪完整构成，仍须补足相应非法占有目的等条件。

$A_I$需按所选解释论证。没有它，不能因受骗人“帮忙把东西递出来”就当然认定三角诈骗既遂；也不能不经论证就把所有第三人交付视为盗窃。

\begin{modelprop}[关系连接必须保留同一主体]
从$D(d,p)$与$P(q,b)$不能直接连接出“$p$受骗后处分$b$财产”的链，除非另有$p=q$或足以连接两人的具体机制。
\end{modelprop}
\begin{proof}
取$p\ne q$，令$p$确实受骗而未处分，令$q$基于无关原因处分。两项原前提都成立，但拟推出的连续机制不存在。因此缺少连接前提时推论不成立。
\end{proof}

即使补入$p=q$，也只解决主体同一问题；仍要证明错误确由有关欺骗引起、处分确由该错误引起。对多人收款、公司付款、自动设备取财等题目，每根箭头写明主体和机制，能够发现把甲的错误认识、乙的交付和丙的损失随意拼接的问题。

\subsection{金额的类型约束}
将犯罪数额、损失、获利、追缴、赔偿分别记为带标签的量
\[
 v_c:\mathsf{CrimeAmount},\quad v_l:\mathsf{Loss},\quad
 v_g:\mathsf{Gain},\quad v_r:\mathsf{Recovery},\quad v_d:\mathsf{Compensation}.
\]
它们可能数值相等，但相等必须有依据。例如退赔使某项未弥补损失减少，不当然使原行为阶段、犯罪数额或既遂事实倒退消失。若某条规范规定退赔影响某种处罚后果，应把该规则单独接入，不能私设恒等式$v_c=v_l-v_d$适用于全部罪名。

\memorycard{一张财产题卡只准使用五栏：“原占有—取得机制—处分主体与地位—新控制—主观目的”。金额另外开五个格，不许用一个“钱数”包办所有判断。}

\section{职务犯罪：身份标签必须连接到具体职务关系}
本节对应第二十三、二十四讲。把“某人有某身份”与“其利用的正是该职务便利”分开。设$O(p,j,t)$表示人在时点$t$任职$j$，$A(j,a)$表示该职务对事项$a$具有相应权限、管理或职责联系，$U(p,j,a)$表示行为人实际利用这条联系。仅有$O$不能推出$U$。

\begin{modelprop}[任职事实不能独立证明职务利用]
从$O(p,j,t)$本身不能推出$U(p,j,a)$。
\end{modelprop}
\begin{proof}
取一个人在该时点任职，但案件事项与其职务无所要求的联系，或者其行为并未利用职务便利的模型。$O$为真而$U$为假，构成反模型。
\end{proof}

这是区分职务便利、工作接近机会、私人关系与影响力路径的基本检查。具体何种联系满足贪污、受贿、斡旋受贿或利用影响力受贿等构成，应分别根据法条展开，不能把“认识官员”当作统一充分条件。\lawref{L1}{第382、385、388、388条之一}

\subsection{中间联系的不同路径不能合并}
以请托人$q$、中间人$m$、实际办理事项者$s$为三类角色。先写$q\to m$的利益给付或请托关系，再写$m\to s$的联系性质，最后写$s\to a$的职务行为。$m\to s$可能属于职权或地位形成的便利，也可能属于近亲属、关系密切者或离职影响力等法定路径；必须保留构成所需身份、时间、便利、利益性质与共同故意条件。只有结构看似同为三个人，不足以把不同罪名视为同一个模型。

\subsection{数额阈值中的析取演算}
2026年解释（二）第一条对单位受贿“情节严重”所列的一组金额路径，在以万元为单位并限定相应适用范围后，可写成
\[
 Q(v,E)=(v\ge20)\vee(10\le v<20\wedge E),
\]
$E$表示具备该条列明的相关情形。\lawref{L8}{第1条} 这是该项门槛的抽取，不是单位受贿罪完整构成。

当$v=15$时，$Q(15,E)=0\vee(1\wedge E)=E$；当$v=20$时，$Q(20,E)=1$，不再需要通过第二金额分支。把这一规则记成“统一20万元起”会遗漏$10\le v<20$且$E=1$的分支；记成“统一10万元起”又会把$E=0$的案件错误纳入。

\memorycard{职务题画三层：人的身份、与事项的职务联系、实际利用方式。金额题画区间，并在较低区间写额外情节变量；边界点10、20各算一次。}

\section{覆盖二十五讲的闭卷演算安排}
这张表把每讲的记忆任务指定为一个可检查的动作。动作完成只说明会重建核心结构，不能替代该讲全部法条和案例学习。

\begingroup\small\linespread{1}\selectfont
\setlength{\LTpre}{4pt}\setlength{\LTpost}{4pt}
\begin{longtable}{p{.09\textwidth}p{.40\textwidth}p{.39\textwidth}}
\toprule
讲次 & 闭卷必须写出的结构 & 用来检查的最小变式\\\midrule
\endfirsthead
\toprule 讲次 & 闭卷必须写出的结构 & 用来检查的最小变式\\\midrule\endhead
1 & 规范依据、解释、空间、时间四层 & 只改变新旧法的轻重关系\\
2 & 完整构成与阻却、责任的依赖图 & 只删去一个必要条件\\
3 & 法益、对象、结果三个不同变量 & 有对象损伤但无所称法益侵害\\
4 & 自然人、身份、单位归责入口 & 员工身份不变而行为归属改变\\
5 & 义务来源、能力、未履行、归责链 & 只改变现实履行能力\\
6 & 基本结果与加重结果分层 & 有较重结果但无所需归责\\
7 & 事实因果与规范归责两张表 & 只改变介入因素的作用机制\\
8 & 认识对象、希望或放任、注意义务 & 只改变对特定结果的认识态度\\
9 & 基础条件、特殊防卫、过当双条件 & 保持损害不变而改变限度条件\\
10 & 年龄区间和各区间附加条件 & 保持罪名不变而跨越年龄边界\\
11 & 阶段、自动性、必要有效性 & 只改变救助是否实际有效\\
12 & 每人一行、每结果一列的矩阵 & 加入一项超出共同故意的行为\\
13 & 行为单元、候选罪名、处理依据 & 同一结果数而行为单元改变\\
14 & 主刑、附加刑和非刑罚措施的类型 & 保留名称近似而更换法律性质\\
15 & 法定刑幅、情节、并罚运算次序 & 总和从34变35；漏罪改为新罪\\
16 & 执行状态和时效起算、停止问题 & 只改变法定最高刑或追诉事件\\
17 & 构成规定、注意规定、拟制入口 & 删除拟制条文后重算候选罪名\\
18 & 法益、行为方式、目的及身份 & 只改变控制他人的目的与机制\\
19 & 占有转换、取得机制、处分关系 & 终点不变而把欺骗改成秘密取走\\
20 & 特定对象风险与公共危险的区别 & 只改变危险可扩散的范围\\
21 & 行政前提与刑事构成两层 & 行政违法成立而严重性条件缺失\\
22 & 具体对象、行为、身份及法定程度 & 只改变诉讼角色或违禁品行为方式\\
23 & 身份、职务联系、利益、交易路径 & 只改变中间人与办理者的联系性质\\
24 & 职责、违背职责、结果、归责 & 有任职身份但无该项具体职责\\
25 & 具体行为、境外关系、对象与主观条件 & 只有境外联系而无所要求行为\\\bottomrule
\end{longtable}
\endgroup

每次检索按三轮进行：第一轮只写变量与规则，不翻答案；第二轮为漏掉的条件造一个反例；第三轮解释该条件为何改变推导。复习记录不只写对错，还应记下失败发生在哪一条边：没有记住规范、没有识别事实、把必要条件当充分条件、误连主体、忽略例外，或者套错版本。这样才能决定下一次应背法条、练案例还是重画推导图。

\modelbound{本补篇证明的对象是明确模型中的逻辑、集合与算术命题。规范是否成立、学说如何选择、事实是否证明，仍须分别论证。前提透明、推导可检查、反例可定位，是这里“严格数学化”的具体含义；不能把它扩张为法律判断已经成为纯计算。}


\chapter{综合演算、反例与复习落实}

\section{一道案件怎样完整作答}
将答题纸固定为六栏：结论候选、规范条件、已知事实、证据或题设来源、仍缺信息、反例变式。复杂案件分人、分行为、分时段填写。先在每一行得出有限结论，再按共犯和罪数规则合并，不从开头就写一个包办全案的罪名。

\caseinput{甲乙约定盗取丙的相机；乙望风。甲入内后发现丙正在持刀严重攻击无关人员丁，甲为救丁将丙制伏。其后甲乙仍按原计划取走相机。甲在离开后主动返还相机并供述。假定各人身份、责任条件及财产犯罪入罪条件已给定。}
\begin{enumerate}
\item 先切分行为：先前计划与帮助、制止侵害、取财、事后返还供述，四段分别记录。
\item 制止侵害段检查正当防卫；甲原有盗窃计划，不当然消灭其在另一行为中可能存在的防卫性质，但必须查清真实目的及防卫条件。
\item 取财段检查盗窃构成与既遂；制止丙的不法侵害并不授权甲取得丙的财物。
\item 对乙分别审查共同故意与望风贡献；不能因甲作出救助行为就将乙的参与自动删除。
\item 事后返还通常不能倒转已经完成的构成事实；是否自首、退赔、从宽，要按各自条件另行分析。
\item 最终分别写明：哪些行为可能阻却违法，哪些行为构成何种候选犯罪，哪些事后事实影响处罚，哪些事实仍需补充。
\end{enumerate}
该演算展示为什么“好人或坏人”的整体印象不能替代行为分析。同一个人对不同对象、在不同时间可能实施性质不同的行为。

\section{十道闭卷检索练习}
先遮住下一节，逐题写出至少三行理由。每题都要求指出适用前提；不以只写“对”“错”计为完成。
\begin{enumerate}
\item 已知$P\Rightarrow Q$且$Q$成立，是否必然有$P$？给一个真值反例，并给一个刑法学习中的对应例子。
\item 案件某罪要求$A\wedge B\wedge C$，现$A=1,B=?,C=1$。工作表应怎样记录？这与“已经证明无罪事实”有什么区别？
\item 一般防卫案件确认造成重大损害，但尚未确认明显超过必要限度。能否仅据此认定防卫过当？
\item 行为时13周岁，故意伤害致重伤，但无死亡、无严重残疾。仅因满12周岁能否按低龄例外追责？
\item 行为人着手后真诚后悔，但没有有效阻止结果发生。能否只凭后悔认定中止？还须区分什么？
\item 盗窃财物数额未达当地通常标准，是否据此一概排除盗窃罪？写出析取结构。
\item 甲与乙共同盗窃，甲另起杀意杀人。能否只凭共同盗窃推出乙负故意杀人共犯责任？
\item 三罪有期徒刑分别为7年、11年、12年，本节通常并罚模型的区间是多少？若分别为10年、12年、15年呢？
\item 某罪法定最高刑为7年，实际判处2年。通常追诉期限分档依据哪个数字？
\item 在第23讲引用的2026年解释金额条件中，$v=15$万元且$E=0$，$Q$为何值？若$E=1$又如何？这能否独立证明单位受贿罪？
\end{enumerate}

\section{逐步解析}
\subsection{第1题：肯定后件的错误}
不必然。取$P=0,Q=1$，则$P\Rightarrow Q$为真，$Q$也为真，但$P$为假。例：某个法定情节可以导致从宽，并不意味着任何从宽结果只能由该情节产生；其他规范路径也可能产生从宽。先列替代路径，再说明不能倒推。

\subsection{第2题：信息缺口}
按本册工作表，$1\wedge?\wedge1=?$。应列出$B$所需事实和证据，不得默认补为$1$。证明要求未达到时如何处理，须适用刑事诉讼规则；逻辑上不能把“尚未证明$B$”改写为“已经证明$\neg B$”。

\subsection{第3题：两项条件的合取}
不能。需要$X\wedge H$，仅知$H=1$而$X$未明，合取尚未得到证明。还须先确认行为具有防卫性质并排除滥用防卫等情形，检查是否适用特殊防卫；不能只在一张脱离前提的真值表上作最终裁判。

\subsection{第4题：年龄规则的限定性}
不能。$12\le a<14$仅进入例外规则的入口。按给定事实，死亡为假，特别残忍手段致重伤造成严重残疾这一组合亦未成立，因此例外条件欠缺。情节恶劣与最高检核准追诉也不能替代该结果条件。

\subsection{第5题：自动性与有效性分离}
不能只凭后悔。先区分是否已经实施终了、是否需要积极防止结果，以及结果是否实际被有效阻止。真诚后悔支持自动性讨论，但不自动证明有效性；结果已发生还须检查归责、故意及其他条件。已经造成的损害和具体犯罪的既遂标准也需独立分析。

\subsection{第6题：析取的一条路径失败}
不能一概排除。记入罪路径为$A\vee M\vee H\vee W\vee P$。$A=0$只表示金额路径失败；若其他一个法定路径成立，整个析取仍可能成立。之后还需核对行为、主观、责任等全部条件。

\subsection{第7题：分别划定故意范围}
不能。先确认乙的共同故意是否包括杀人及其贡献，再处理甲的实行过限。共同盗窃只是对盗窃部分的关联事实，不能无条件推出杀人共同故意。是否有其他责任仍需额外事实。

\subsection{第8题：先算总和，再套上限}
第一组$m=12,S=30<35$，所以$U=20$，得到$12\le y\le20$。第二组$m=15,S=37\ge35$，所以$U=25$，得到$15\le y\le25$。这是通常有期徒刑并罚的法定区间检验；具体执行刑期仍需依法酌定，异种刑罚或特别规定不在本算式范围内。

\subsection{第9题：分档变量的类型}
应以该罪适用的法定最高刑7年分档，而非宣告刑2年。由$5\le7<10$，得到通常期限10年；再检查起算、期间内再犯罪、不受期限限制等条件。只知道7年并不足以得到具体的日历截止日期。

\subsection{第10题：分支与整体构成}
$15<20$，且$10\le15<20$。当$E=0$时，两条路径均不成立，$Q=0$；当$E=1$时，第二路径成立，$Q=1$。$Q$只代表该解释的相应情节条件，不能独立证明主体、受贿行为、主观等其余构成。

\section{怎样把记忆任务落实到每一讲}
每天选择三至五张卡即可，数量按掌握情况调整。第一遍写条件；第二遍给反例；第三遍只改变一个事实重新作答。复习时把错误直接定位到模型，不笼统写“粗心”。例如把数额换成损失，属于变量类型错误；把“可以”写成“应当”，属于规范模态错误；把事后知道写成事前明知，属于时间错误；把旧解释数字用于新情形，属于版本错误。

\begin{center}
\begin{tabular}{p{.16\textwidth}p{.33\textwidth}p{.40\textwidth}}
\toprule
错误类型 & 自己写下的错误 & 具体修正动作\\\midrule
漏条件 & 有死亡就认定过当 & 补写明显超限，并复核防卫基础条件\\
错连接词 & 不能辨认且不能控制 & 对照法条恢复“或者”，做真值反例\\
错对象 & 金融交易销售额代替所得 & 给每个金额添加口径与来源\\
错时间 & 事后后悔等于事前无故意 & 画行为前、行为中、行为后三段轴\\
错版本 & 沿用旧单位受贿门槛 & 核对2026年解释适用时点及具体条件\\\bottomrule
\end{tabular}
\end{center}

\chapter{覆盖范围、规范出处与后续校订边界}

\section{本册已经做了什么}
本册按底本目录逐一设置二十五讲的对应节，完成核心概念的输入结构、法律前提与模型推导的区分、每讲教学设例或变式、记忆检索方法，并对重点总论与财产犯罪给出较详细演算。模型命题的证明限于明确前提下的逻辑或算术结论；没有把法律学说、作者评价或未知事实伪装成数学定理。

本册\textbf{不是原书全部案例的逐条形式化替代版本}。尤其是分论数百个具体罪名、全部司法解释、地方金额标准及历史真题的逐题争议，仍需建立逐条对照表再核验。本册覆盖某一讲，表示已经提供该讲的核心学习模型，不表示该讲的每个细节已经逐一校订。没有完成Lean或其他证明助手的机器验证；排版编译成功也不等于法律结论经过形式验证。

\section{二十五讲的模型对应表}
\begingroup\small\linespread{1}\selectfont
\setlength{\LTpre}{4pt}\setlength{\LTpost}{4pt}
\begin{longtable}{p{.10\textwidth}p{.23\textwidth}p{.52\textwidth}}
\toprule
讲次 & 原书主题 & 本册核心检查对象\\\midrule
\endhead
1 & 刑法论 & 规范依据、解释集合、时间与地域连接\\
2 & 犯罪构成 & 条件合取、三段论、必要与充分条件\\
3 & 犯罪客体 & 对象与法益的二部关系\\
4 & 犯罪主体 & 自然人、身份、单位归属\\
5 & 行为 & 不作为的义务、能力、不履行与等价\\
6 & 结果 & 实害、危险、结果加重的分层\\
7 & 因果关系 & 机制、介入因素、竞争解释\\
8 & 主观要件 & 特定结果的认识态度、现实与认识二栏\\
9 & 排除违法事由 & 防卫基础、双条件过当、避险与承诺\\
10 & 排除责任事由 & 年龄分段、责任能力与理论边界\\
11 & 犯罪形态 & 着手、既遂、停止原因与有效性\\
12 & 共同犯罪 & 个人贡献、故意范围、时间与关系图\\
13 & 罪数 & 行为—罪名矩阵与规范处理关系\\
14 & 刑罚体系 & 刑罚类型与参数，避免异类加法\\
15 & 刑罚裁量 & 模态、区间、并罚、漏罪与新罪\\
16 & 执行和消灭 & 刑罚状态、假释、时效分段与例外\\
17 & 分论概说 & 罪状字段、援引规则与法律拟制\\
18 & 人身犯罪 & 不同法益、主观、结果与特别规则\\
19 & 财产犯罪 & 占有、取财机制、处分、金额与阶段\\
20 & 公共安全 & 危险范围、方法、六类犯罪检索\\
21 & 经济秩序 & 行政与刑事条件、金额口径、版本\\
22 & 社会管理 & 九组对象与行为、鉴定和主观条件\\
23 & 贪污贿赂 & 身份职务财产三线，2026年规则样例\\
24 & 渎职 & 职责、违反、结果归责与主观\\
25 & 国家安全 & 具体罪状、对象性质与证据关系\\\bottomrule
\end{longtable}
\endgroup

\section{规范来源及核验记录}
以下列出实际用于本册核心规则核对的公开规范文本。核验日期为2026年11月20日。一般刑法条文以国家机关公布的修正案（十一）整理文本为基础，并对相关修正案（十二）条款另行核对；不能把该基础页面误称为已经包含所有后续修订。关于2026年贪污贿赂新解释，另列独立来源。来源表不代表已经穷尽全国所有法律、司法解释与地方规定。

\begingroup\small
\begin{description}[itemsep=4pt,parsep=0pt,topsep=4pt]
\item[\hypertarget{L1}{L1}] 《中华人民共和国刑法》，国家信访局公布，来源注明国家法律法规数据库，整理到刑法修正案（十一）。用于核对总则及相关分则的基础条文；修正案（十二）另见L3。\par\url{https://www.gjxfj.gov.cn/2020-12/26/c_139953427.htm}
\item[\hypertarget{L2}{L2}] 《中华人民共和国刑法修正案（十一）》，2020年12月26日通过，2021年3月1日起施行，最高人民检察院公开文本；本册重点核对第17条年龄规定及照护职责有关新条文。\par\url{https://www.spp.gov.cn/spp/fl/202012/t20201227_503700.shtml}
\item[\hypertarget{L3}{L3}] 《中华人民共和国刑法修正案（十二）》，2023年12月29日通过，2024年3月1日起施行，新华社受权发布全文；用于核对相关公司企业人员犯罪和行贿条款修订。\par\url{https://www.news.cn/politics/20231229/387cc35b798f465880139ad212b0b8c8/c.html}
\item[\hypertarget{L4}{L4}] 最高人民法院、最高人民检察院、公安部《关于依法适用正当防卫制度的指导意见》，法发〔2020〕31号，2020年8月28日印发。\par\url{https://www.court.gov.cn/zixun/xiangqing/251611.html}
\item[\hypertarget{L5}{L5}] 《中华人民共和国刑事诉讼法》，市场监管总局公布的2018年修正文本；本册核对第55条证据标准。\par\url{https://www.samr.gov.cn/jjz/djfg/gjflfg/art/2023/art_929f3de7d593441bbe63fe0eecd05e10.html}
\item[\hypertarget{L6}{L6}] 《中华人民共和国刑法》第69至71条数罪并罚、第99条有关用语解释，见L1同一官方整理文本；本册只将符合限定条件的有期徒刑并罚改写为区间。\par\url{https://www.gjxfj.gov.cn/2020-12/26/c_139953427.htm}
\item[\hypertarget{L7}{L7}] 最高人民法院、最高人民检察院《关于办理盗窃刑事案件适用法律若干问题的解释》，法释〔2013〕8号，2013年4月4日起施行；本册核对数额标准框架、多次盗窃等路径与未遂规定。\par\url{https://www.court.gov.cn/zixun/xiangqing/5228.html}
\item[\hypertarget{L8}{L8}] 最高人民法院、最高人民检察院《关于办理贪污贿赂刑事案件适用法律若干问题的解释（二）》，法释〔2026〕6号，2026年5月1日起施行；本册具体演算第1条中的金额与情节结构，并定位第8、11至17条有关变化。\par\url{https://www.court.gov.cn/fabu/xiangqing/497181.html}
\item[\hypertarget{L9}{L9}] 《中华人民共和国行政许可法》，市场监管总局公布的2019年修正文本，第54条；用于核查“不得指定教材”与“官方编写出版”的逻辑区别。\par\href{https://www.samr.gov.cn/zw/zfxxgk/fdzdgknr/tzsbs/art/2023/art_3f97055b34d241dbaba5d4a81757901f.html}{市场监管总局官网：《中华人民共和国行政许可法》全文}
\end{description}
\endgroup

\section{继续逐条校订时的合格记录}
每条原书规则均应保存：原句、对应条文或解释、行为时间、适用版本、形式化表达、采用的争议立场、至少一个正例、至少一个反例，以及实际通过的检查。原文若只有单向规则，就保留单向箭头；要升级为双向等价，必须补证逆向。原文若是经验判断，就注明经验性质；要写成普遍命题，需要足以支持全称断言的前提。

若发现一条规则与新解释不一致，应记录具体冲突，而不是为维持整齐统一而悄悄删除原说。若案例事实不足，应列出待补问题；若两个解释都可能成立，应分别给出条件化结论。严格性体现为每个结论能追到前提，以及每个前提能追到来源。
